\SourcePage{402}
\SGAsection{10}{群概形及带算子群的概形中的投射极限}
\textbf{10.0。}回顾 EGA IV$_3$，§8.8 的基本结果。给定如下情形：$S_0$ 是\Term{quasi-compact}{拟紧}且\Term{quasi-separated}{拟分离}的\Term{scheme}{概形}，$I$ 是递增\Term{filtered-category}{滤过的}\Term{preordered-set}{预序集}，$(\mathcal A_i)_{i\in I}$ 是交换\Term{quasi-coherent}{拟凝聚} $\mathcal O_{S_0}$-\Term{algebra}{代数}的\Term{inductive-system}{归纳系统}，$\mathcal A=\varinjlim\mathcal A_i$，对 $i\in I$ 置 $S_i=\operatorname{Spec}\mathcal A_i$，并置 $S=\operatorname{Spec}\mathcal A$。\footnote{编者注：注意，$S$ 在 $S_0$ 上是\Term{affine}{仿射}的，因而\Term{quasi-compact}{拟紧}且\Term{quasi-separated}{拟分离}；参见下文编者注（92）。}那么，$S$ 上\Term{finite-presentation}{有限表示}的\Term{scheme}{概形}所成的\Term{category}{范畴}，在\Term{equivalence}{等价}意义下由各 $S_i$ 上\Term{finite-presentation}{有限表示}的\Term{scheme}{概形}所成的\Term{category}{范畴}、这些\Term{category}{范畴}之间的\Term{functor}{函子}
\[
\rho_{ji}:X_i\longmapsto X_i\times_{S_i}S_j\quad(i\leq j),
\]
以及\Term{transitivity-isomorphism}{传递性同构} $\rho_{kj}\circ\rho_{ji}\xrightarrow{\sim}\rho_{ki}$ 决定。

\SourcePage{403}
具体来说，给定 $j\in I$ 及\Term{finite-presentation}{有限表示}的 $S_j$-\Term{scheme}{概形} $X_j$，对每个 $i\geq j$ 置 $X_i=X_j\times_{S_j}S_i$，并置 $X=X_j\times_{S_j}S$。则有（EGA IV$_3$，8.8.2）：

（i）给定 $j\in I$ 及两个\Term{finite-presentation}{有限表示}的 $S_j$-\Term{scheme}{概形} $X_j,Y_j$，\Term{canonical}{典范}映射
\[
\varinjlim_{i\geq j}\operatorname{Hom}_{S_i}(X_i,Y_i)\longrightarrow\operatorname{Hom}_S(X,Y)
\]
是\Term{bijection}{双射}。

（ii）对每个\Term{finite-presentation}{有限表示}的 $S$-\Term{scheme}{概形} $X$，存在指标 $j\in I$、\Term{finite-presentation}{有限表示}的 $S_j$-\Term{scheme}{概形} $X_j$ 及 $S$-\Term{isomorphism}{同构} $X\xrightarrow{\sim}X_j\times_{S_j}S$。

由此\Term{pushout}{推出}（EGA IV$_3$，8.8.3），只要 $D$ 是\Term{finite-presentation}{有限表示}的 $S$-\Term{scheme}{概形}\Term{category}{范畴}中只涉及\Term{finite}{有限}多个\Term{object}{对象}和\Term{arrow}{箭头}的图，就能找到指标 $i\in I$ 及\Term{finite-presentation}{有限表示}的 $S_i$-\Term{scheme}{概形}\Term{category}{范畴}中的图 $D_i$，使 $D$ 在\Term{isomorphism}{同构}意义下由 $D_i$ 经\Term{base-change}{基变换} $S\to S_i$ 得到。甚至可以选取 $i,D_i$，使 $D$ 的每个\Term{cartesian}{笛卡尔}方块都来自 $D_i$ 的一个\Term{cartesian}{笛卡尔}方块。

\textbf{10.1。}此外，\Term{morphism}{态射}的许多常见性质在\Term{base-change}{基变换}下\Term{stable}{稳定}，并具有如下性质：
\begin{quote}
设 $u:X\to Y$ 是\Term{finite-presentation}{有限表示}的 $S$-\Term{scheme}{概形}之间的 $S$-\Term{morphism}{态射}。由 10.0，它由\Term{finite-presentation}{有限表示}的 $S_j$-\Term{scheme}{概形}之间的某个 $S_j$-\Term{morphism}{态射} $u_j:X_j\to Y_j$ 经\Term{base-change}{基变换}得到。那么，$u$ 具有性质 $P$ 的充要条件是，存在 $i\geq j$，使 $u_i=u_j\times_{S_j}S_i$ 具有性质 $P$。
\end{quote}
当 $P$ 是下列\Term{morphism}{态射}性质之一时，上述断言成立：\Term{separated}{分离}、\Term{surjectivity}{满射}、\Term{radicial}{根性}、\Term{affine}{仿射}、\Term{quasi-affine}{拟仿射}、\Term{finite}{有限}、\Term{quasi-finite}{拟有限}、\Term{proper}{固有}、\Term{projective}{射影}、\Term{quasi-projective}{拟射影}、\Term{isomorphism}{同构}、\Term{monomorphism}{单态射}、\Term{immersion}{浸入}、\Term{open-immersion}{开浸入}、\Term{closed-immersion}{闭浸入}（EGA IV$_3$，8.10.5）、\Term{flat}{平坦}（EGA IV$_3$，11.2.6）、\Term{smooth}{光滑}、\Term{unramified}{非分歧}或\Term{etale}{平展}（EGA IV$_4$，17.7.8）。\footnote{编者注：补入了“\Term{flat}{平坦}”一词，并将 17.7.6 改为 17.7.8。}

注意，当 $P$ 是相对于 $(\mathrm{fppf})$ \Term{topology}{拓扑}的\Term{covering}{覆盖}性质时亦然。事实上，给定两个\Term{finite-presentation}{有限表示}的 $S$-\Term{scheme}{概形} $X,Y$ 及 $S$-\Term{morphism}{态射} $u:X\to Y$，由第 IV 讲 6.3.1（i）\footnote{编者注：还要用到 $Y$ 在 $S$ 上\Term{finite-presentation}{有限表示}，因而\Term{quasi-compact}{拟紧}这一事实。}可知，$u$ 相对于 $(\mathrm{fppf})$ \Term{topology}{拓扑}为\Term{covering}{覆盖}的充要条件是，存在一个 $S$-\Term{scheme}{概形} $Z$ 及\Term{faithfully-flat}{忠实平坦}且\Term{finite-presentation}{有限表示}的 $S$-\Term{morphism}{态射} $v:Z\to Y$，使其经 $u$ \Term{resolution}{分解}。

本第 10 节的目的是，对\Term{finite-presentation}{有限表示}的 $S$-\Term{group-scheme}{群概形}的\Term{category}{范畴}、带\Term{operator-group}{算子群}的 $S$-\Term{scheme}{概形}的\Term{category}{范畴}，以及群的\Term{monomorphism}{单态射}的某些性质（\Term{normal-subgroup}{正规}、\Term{central}{中心}且\Term{quotient-sheaf}{商层}\Term{representable}{可表}等），给出这类结果的变式。

这一类型的两个初步结果如下。（以下 10.2 至 10.9 沿用 10.0 引入的记号。）

\SourcePage{404}
\begin{SGAresult}{\Term{lemma}{引理}}{10.2}
设 $G_j,H_j$ 是两个\Term{finite-presentation}{有限表示}的 $S_j$-\Term{group-scheme}{群概形}；对每个 $i\geq j$，置 $G_i=G_j\times_{S_j}S_i$，$G=G_j\times_{S_j}S$，并类似地定义 $H_i,H$。则下列\Term{canonical}{典范}映射是\Term{bijection}{双射}：
\[
\varinjlim_{i\geq j}\operatorname{Hom}_{S_i\text{-gr.}}(G_i,H_i)\longrightarrow\operatorname{Hom}_{S\text{-gr.}}(G,H).
\]
\end{SGAresult}
\begin{SGAresult}{\Term{lemma}{引理}}{10.3}
设 $G$ 是\Term{finite-presentation}{有限表示}的 $S$-\Term{group-scheme}{群概形}。则存在 $j\in I$、\Term{finite-presentation}{有限表示}的 $S_j$-\Term{group-scheme}{群概形} $G_j$，以及 $S$-\Term{group-scheme}{群概形}的\Term{isomorphism}{同构} $G\simeq G_j\times_{S_j}S$。
\end{SGAresult}
结合 EGA $0_{\mathrm{III}}$，8.2.5、8.2.6 对 $S$-\Term{group-structure}{群结构}的解释，\footnote{编者注：即用 $S$-\Term{morphism}{态射}的\Term{commutative-diagram}{交换图}来表述。}10.2、10.3 容易由 10.0、10.1 得到。

\begin{SGAresult}{\Term{lemma}{引理}}{10.4}
设 $u:G\to H$ 是\Term{finite-presentation}{有限表示}的 $S$-\Term{group-scheme}{群概形}之间的 $S$-\Term{group-morphism}{群态射}。由 10.3、10.2，$u$ 由\Term{finite-presentation}{有限表示}的 $S_j$-\Term{group-scheme}{群概形}之间的\Term{morphism}{态射} $u_j:G_j\to H_j$ 经\Term{base-change}{基变换}得到。那么，$u$ 为\Term{central}{中心}\Term{monomorphism}{单态射}（相应地，为\Term{normal-subgroup}{正规}\Term{monomorphism}{单态射}）的充要条件是，存在 $i\geq j$，使 $u_i=u_j\times_{S_j}S_i$ 具有同一性质。
\end{SGAresult}
结合 6.7 给出的群的\Term{central}{中心}或\Term{normal-subgroup}{正规}\Term{monomorphism}{单态射}的刻画，这立即由 10.0、10.1 得到。

\begin{SGAresult}{推论}{10.5}
设 $G_j$ 是\Term{finite-presentation}{有限表示}的 $S_j$-\Term{group-scheme}{群概形}。$G_j\times_{S_j}S$ 交换的充要条件是，存在 $i\geq j$，使 $G_j\times_{S_j}S_i$ 交换。
\end{SGAresult}
事实上，一个 $S$-\Term{group-scheme}{群概形}交换，\Term{equivalence}{等价}于把它看作自身的\Term{subgroup-scheme}{子群概形}时它是\Term{central}{中心}的。

\begin{SGAresult}{命题}{10.6}
设 $G_j$ 是\Term{finite-presentation}{有限表示}的 $S_j$-\Term{group-scheme}{群概形}，$G'_j$ 是 $G_j$ 的\Term{subgroup-scheme}{子群概形}，且在 $S_j$ 上\Term{flat}{平坦}、\Term{finite-presentation}{有限表示}。则 $(G_j\times_{S_j}S)/(G'_j\times_{S_j}S)$ 相对于 $(\mathrm{fpqc})$ \Term{topology}{拓扑}\Term{representable}{可表}的充要条件是，存在 $i\geq j$，使 $(G_j\times_{S_j}S_i)/(G'_j\times_{S_j}S_i)$ 亦然。
\end{SGAresult}
这是下列更一般的\Term{lemma}{引理}的推论。
\begin{SGAresult}{\Term{lemma}{引理}}{10.7}
设 $X_j$ 是\Term{finite-presentation}{有限表示}的 $S_j$-\Term{scheme}{概形}，$R_j$ 是 $X_j$ 上\Term{flat}{平坦}且\Term{finite-presentation}{有限表示}的\Term{equivalence-relation}{等价关系}。\footnote{编者注：意即复合 $R_j\hookrightarrow X_j\times_{S_j}X_j\xrightarrow{\operatorname{pr}_1}X_j$ 是\Term{flat}{平坦}且\Term{finite-presentation}{有限表示}的。}那么，相对于\Term{topology}{拓扑} $T=(\mathrm{fppf})$ 或 $(\mathrm{fpqc})$ 的\Term{quotient-sheaf}{商层} $(X_j\times_{S_j}S)/(R_j\times_{S_j}S)$ \Term{representable}{可表}的充要条件是，存在 $i\geq j$，使相对于\Term{topology}{拓扑} $T$ 的\Term{quotient-sheaf}{商层} $(X_j\times_{S_j}S_i)/(R_j\times_{S_j}S_i)$ \Term{representable}{可表}。
\end{SGAresult}
结合 10.0 中回顾的 EGA IV$_2$，8.8.2、8.8.3、8.10.5、11.2.6 的陈述，本\Term{lemma}{引理}由下列结果\Term{pushout}{推出}。
\begin{SGAresult}{\Term{lemma}{引理}}{10.8}
设 $T$ 为 $(\mathrm{fppf})$ 或 $(\mathrm{fpqc})$ \Term{topology}{拓扑}，$X$ 为\Term{finite-presentation}{有限表示}（相应地，\Term{locally-finite-presentation}{局部有限表示}）的 $S$-\Term{scheme}{概形}，$R$ 为 $X$ 上由\Term{monomorphism}{单态射} $v:R\to X\times_S X$ 定义的\Term{equivalence-relation}{等价关系}，使 $\operatorname{pr}_1\circ v$ 为\Term{flat}{平坦}且\Term{finite-presentation}{有限表示}（相应地，为\Term{flat}{平坦}且\Term{locally-finite-presentation}{局部有限表示}）。
\SourcePage{405}
则下列条件\Term{equivalence}{等价}：

（i）相对于\Term{topology}{拓扑} $T$ 的\Term{quotient-sheaf}{商层} $X/R$ \Term{representable}{可表}。

（ii）存在\Term{finite-presentation}{有限表示}（相应地，\Term{locally-finite-presentation}{局部有限表示}）的 $S$-\Term{scheme}{概形} $Y$ 及\Term{faithfully-flat}{忠实平坦}\Term{morphism}{态射} $p:X\to Y$，使下图为\Term{cartesian}{笛卡尔}图：
\[
\begin{tikzcd}
R\arrow[r,"\operatorname{pr}_1\circ v"]\arrow[d,"\operatorname{pr}_2\circ v"']&X\arrow[d,"p"]\\
X\arrow[r,"p"']&Y.
\end{tikzcd}\tag{D}
\]
\end{SGAresult}
\begin{proof}
先注意，由第 IV 讲 3.3.2、4.4.3，相对于\Term{topology}{拓扑} $T$ 的\Term{sheaf}{层} $X/R$ 由 $Y$ 表示的充要条件是，图（D）为\Term{cartesian}{笛卡尔}图，且 $p$ 相对于\Term{topology}{拓扑} $T$ 为\Term{covering}{覆盖}。

证明（i）\Term{pushout}{推出}（ii）。假设（i）蕴含图（D）为\Term{cartesian}{笛卡尔}图，因此 $\operatorname{pr}_1\circ v$ 由 $p$ 沿 $p$ 作\Term{base-change}{基变换}得到，且 $p$ 相对于 $(\mathrm{fpqc})$ \Term{topology}{拓扑}为\Term{covering}{覆盖}。于是由 $(\mathrm{fpqc})$ \Term{descent}{下降}（EGA IV$_2$，2.7.1），既然 $\operatorname{pr}_1\circ v$ 是\Term{faithfully-flat}{忠实平坦}且（\Term{local-scheme}{局部}）\Term{finite-presentation}{有限表示}的，$p$ 亦然。又由 EGA IV$_4$，17.7.5，既然 $X$ 在 $S$ 上（\Term{local-scheme}{局部}）\Term{finite-presentation}{有限表示}，$Y$ 亦然。

证明（ii）\Term{pushout}{推出}（i）。只须证明 $p$ 相对于 $(\mathrm{fppf})$ \Term{topology}{拓扑}为\Term{covering}{覆盖}。由假设，$p$ \Term{faithfully-flat}{忠实平坦}；又因 $X,Y$ 在 $S$ 上均\Term{locally-finite-presentation}{局部有限表示}，$p$ 也\Term{locally-finite-presentation}{局部有限表示}（EGA IV$_1$，1.4.3（v））。
\end{proof}

\begin{SGAresult}{\Term{lemma}{引理}}{10.9}
设 $G_j$ 是\Term{finite-presentation}{有限表示}的 $S_j$-\Term{group-scheme}{群概形}，$G=G_j\times_{S_j}S$。则 $G^0$ \Term{representable}{可表}的充要条件是，存在 $i\geq j$，使 $(G_i)^0=(G_j\times_{S_j}S_i)^0$ \Term{representable}{可表}。
\end{SGAresult}
\begin{proof}
此条件充分，因为由 3.3，\Term{functor}{函子} $G\mapsto G^0$ 与\Term{base-change}{基变换}交换。

反之，设 $G^0$ \Term{representable}{可表}。由 3.9，$G^0$ 在 $G$ 中为开，且在 $S$ 上\Term{quasi-compact}{拟紧}；因 $G$ \Term{finite-presentation}{有限表示}，$G^0$ 在 $S$ 上也\Term{finite-presentation}{有限表示}。于是由 10.3、10.1，存在 $i\geq j$ 及 $G_i$ 的开\Term{subgroup-scheme}{子群概形} $H_i$，使 $H_i\times_{S_i}S=G^0$。\Term{structural-morphism}{结构态射} $G^0\to S$ 是\Term{connected}{连通}的，即其\Term{fiber}{纤维}\Term{geometrically-connected}{几何连通}（第 VI$_A$ 讲 2.1.1）。故由 EGA IV$_3$，9.3.3、9.7.7，增大 $i$ 后可设\Term{structural-morphism}{结构态射} $H_i\to S$ 是\Term{connected}{连通}的。于是由 3.10.1，$H_i$ 的底\Term{space}{空间}就是 $(G_i)^0$，因而 $H_i$ 表示 $(G_i)^0$。
\end{proof}

\textbf{10.10。}回顾 10.0 所述情形的两个非常有用的特例（参见 EGA IV$_3$，8.1.2 a）、c））：
\SourcePage{406}

 a）给定\Term{scheme}{概形} $X$ 的一点 $x$，置 $S_0=\operatorname{Spec}\mathbb Z$，并考虑 $x$ 的\Term{affine}{仿射}\Term{open-neighborhood}{开邻域}构成的递减\Term{filtered}{滤过}\Term{projective-system}{投射系统} $(S_i)_{i\in I}$；则 $S=\operatorname{Spec}\mathcal O_{X,x}$。特别地，若 $x$ 为\Term{integral}{整}\Term{scheme}{概形} $X$ 的\Term{generic-point}{泛点}，就得到 $S=\operatorname{Spec}\kappa(x)$。

 b）置 $S_0=\operatorname{Spec}\mathbb Z$，并考虑\Term{affine}{仿射}\Term{scheme}{概形} $S$ 的环中所有\Term{finite-type}{有限型} $\mathbb Z$-\Term{subalgebra}{子代数}构成的族 $(A_i)_{i\in I}$，以包含关系赋予\Term{preorder}{预序}。由于 $A_i$ 都是\Term{noetherian}{诺特}环，这在许多情形下允许从\Term{noetherian}{诺特}情形过渡到一般情形。

现在给出关于 a）中所考察特例的两个结果。\footnote{编者注：注意，证明还要用到 b）的情形。}
\begin{SGAresult}{命题}{10.11}
\footnote{编者注：简化了陈述，并在推论 10.11.1 中单独处理了\Term{subgroup}{子群}的情形。}设 $S$ 为\Term{integral}{整}\Term{scheme}{概形}，$\eta$ 为其\Term{generic-point}{泛点}，$G$ 为\Term{finite-presentation}{有限表示}的 $S$-\Term{group-scheme}{群概形}，$Y,Z$ 为\Term{finite-presentation}{有限表示}的 $S$-\Term{scheme}{概形}，$u,v,i:Y\to G$ 及 $j:Z\to G$ 为 $S$-\Term{scheme}{概形}\Term{morphism}{态射}。假设 $i_\eta:Y_\eta\to G_\eta$ 及 $j_\eta:Z_\eta\to G_\eta$ 都是\Term{closed-immersion}{闭浸入}。

则存在 $S$ 的\Term{nonempty}{非空}\Term{open-subset}{开集} $S'$，使\Term{base-change}{基变换}后得到的\Term{morphism}{态射} $i':Y'\to G'$、$j':Z'\to G'$ 为\Term{closed-immersion}{闭浸入}，且\Term{functor}{函子}
\[
\operatorname{Transp}(u',v'),\quad
\operatorname{Transp}_{G'}(i'(Y'),j'(Z')),\quad
\operatorname{Transpstr}_{G'}(i'(Y'),j'(Z'))
\]
以及相应的
\[
\operatorname{Centr}(u'),\qquad\operatorname{Norm}_{G'}i'(Y')
\]
可由 $G'$ 中在 $S'$ 上\Term{finite-presentation}{有限表示}的闭 $S'$-\Term{subscheme}{子概形}（相应地，闭 $S'$-\Term{subgroup-scheme}{子群概形}）表示。
\end{SGAresult}
\begin{proof}
将先在 10.10 a）的情形中、再在 10.10 b）的情形中应用 10.1 的结果。因为 $G_\eta,Y_\eta,Z_\eta$ 在域 $\kappa(\eta)$ 上\Term{flat}{平坦}，$G_\eta$ 在 $\kappa(\eta)$ 上\Term{separated}{分离}（第 VI$_A$ 讲 0.3），且 $i_\eta,j_\eta$ 为\Term{closed-immersion}{闭浸入}，所以由 10.1，存在 $S$ 的\Term{affine}{仿射}\Term{open-subset}{开集} $S'=\operatorname{Spec}(A')$、$A'$ 的\Term{noetherian}{诺特}\Term{subring}{子环} $A$、在 $A$ 上\Term{flat}{平坦}且\Term{finite-presentation}{有限表示}的 $A$-\Term{scheme}{概形} $G_A,Y_A,Z_A$，以及\Term{morphism}{态射} $u_A,v_A,i_A:Y_A\to G_A$、$j_A:Z_A\to G_A$，使 $G_A$ 是在 $A$ 上\Term{separated}{分离}的 $A$-\Term{group-scheme}{群概形}，$i_A,j_A$ 是\Term{closed-immersion}{闭浸入}，且 $G\times_S S'=G_A\otimes A'$，等等。因在 $S'$ 上考虑的\Term{functor}{函子}由 $\operatorname{Spec}(A)$ 上的对应\Term{functor}{函子}经\Term{base-change}{基变换}得到，只须对后者证明结论。

由 EGA IV$_2$，6.9.2，将 $A$ 换成某个\Term{localization}{局部化} $A_f$（相应地，将 $S'$ 换成\Term{affine}{仿射}\Term{open-subset}{开集} $S'_f$）后，可设 $G_A,Y_A,Z_A$ 在 $A$ 上\Term{essentially-free}{本质自由}（按 6.2.1 的意义）。\footnote{编者注：事实上，$G_A$ 在 $A$ 上\Term{flat}{平坦}且\Term{finite-presentation}{有限表示}，故可用\Term{affine}{仿射}\Term{open-subset}{开集} $G_1,\ldots,G_n$ \Term{covering}{覆盖}，使每个 $\mathcal O(G_i)$ 为\Term{flat}{平坦}且\Term{finite-presentation}{有限表示}的 $A$-\Term{algebra}{代数}。于是由 EGA IV$_2$，6.9.2，存在 $f_i\in A$，使 $\mathcal O(G_i)_{f_i}$ 为 $A_{f_i}$ 上的\Term{free-module}{自由模}；可将 $\operatorname{Spec}(A)$ 换成\Term{affine}{仿射}\Term{open-subset}{开集} $D(f)$，其中 $f=f_1\cdots f_n$；对 $Y_A,Z_A$ 同样处理。}于是由 6.2.4 b）、e），在命题的假设下，所考察的\Term{functor}{函子}可由 $G_A$ 的闭 $A$-\Term{subscheme}{子概形}表示（因 $A$ 是\Term{noetherian}{诺特}环，且 $G_A$ 在 $A$ 上\Term{finite-presentation}{有限表示}，它们在 $A$ 上也\Term{finite-presentation}{有限表示}）；
\SourcePage{407}
在 $\operatorname{Centr}(u_A)$ 及 $\operatorname{Norm}_{G_A}i_A(Y_A)$ 的情形，它们是 $G_A$ 的 $A$-\Term{subgroup-scheme}{子群概形}。
\end{proof}

\begin{SGAresult}{推论}{10.11.1}
设 $S$ 为\Term{integral}{整}\Term{scheme}{概形}，$\eta$ 为其\Term{generic-point}{泛点}，$G,H,K$ 为\Term{finite-presentation}{有限表示}的 $S$-\Term{group-scheme}{群概形}，$i:H\to G$、$j:K\to G$ 为 $S$-群的两个\Term{quasi-compact}{拟紧}\Term{monomorphism}{单态射}。则存在 $S$ 的\Term{nonempty}{非空}\Term{open-subset}{开集} $S'$，使\Term{base-change}{基变换}后得到的\Term{morphism}{态射} $i':H'\to G'$、$j':K'\to G'$ 为\Term{closed-immersion}{闭浸入}，且\Term{functor}{函子}
\[
\operatorname{Transp}_{G'}(H',K'),\quad\operatorname{Transpstr}_{G'}(H',K')
\quad\text{（相应地，}\operatorname{Centr}_{G'}H',\quad\operatorname{Norm}_{G'}H'\text{）}
\]
可由 $G'$ 中在 $S'$ 上\Term{finite-presentation}{有限表示}的闭 $S'$-\Term{subscheme}{子概形}（相应地，闭 $S'$-\Term{subgroup-scheme}{子群概形}）表示。
\end{SGAresult}
由 1.4.2，假设蕴含 $i_\eta,j_\eta$ 为\Term{closed-immersion}{闭浸入}，因而此结论由上一命题\Term{pushout}{推出}。

\begin{SGAresult}{命题}{10.12}
设 $S$ 为\Term{integral}{整}\Term{scheme}{概形}，$G$ 为\Term{finite-presentation}{有限表示}的 $S$-\Term{group-scheme}{群概形}，$A,B$ 为 $G$ 的\Term{subgroup-scheme}{子群概形}，在 $S$ 上\Term{finite-presentation}{有限表示}且\Term{generic-fibre}{泛纤维}\Term{smooth}{光滑}。再假设满足下列条件之一：

 a）$A$ 的\Term{generic-fibre}{泛纤维}\Term{connected}{连通}；

 b）$A,B$ 都是 $G$ 的\Term{normal-subgroup}{正规子群}。

则存在 $S$ 的\Term{nonempty}{非空}\Term{open-subset}{开集} $S'$ 及 $G'=G|_{S'}$ 的闭\Term{subgroup-scheme}{子群概形} $D'$，在 $S'$ 上\Term{finite-presentation}{有限表示}、\Term{fiber}{纤维}\Term{smooth}{光滑}，表示 $G'$ 中 $A'=A|_{S'}$ 和 $B'=B|_{S'}$ 的\Term{commutator}{交换子}\Term{group-presheaf}{群预层}所伴的 $(\mathrm{fppf})$ \Term{sheaf}{层}。在（a）中，$D'$ 的\Term{fiber}{纤维}\Term{connected}{连通}；在（b）中，$D'$ 是 $G'$ 的\Term{normal-subgroup}{正规子群}。

特别地，对每个 $s\in S'$，按 7.2.2 的记号，有 $D'_s=(A_s,B_s)$。
\end{SGAresult}
\begin{proof}
设 $\eta$ 是 $S$ 的\Term{generic-point}{泛点}，置 $H_\eta=(A_\eta,B_\eta)$。由于 $A_\eta,B_\eta$ 都\Term{smooth}{光滑}，在（a）中由 7.8，在（b）中由 7.3（v），$H_\eta$ 是\Term{connected}{连通}的（相应地，是 $G_\eta$ 的\Term{normal-subgroup}{正规子群}）。

我们处在与 10.10（a）相应的 10.0 的情形。因此由 10.3、10.1，存在 $S$ 的\Term{nonempty}{非空}\Term{open-subset}{开集} $S'$ 及在 $S'$ 上\Term{finite-presentation}{有限表示}、在 $G'$ 中闭的 $S'$-\Term{subgroup-scheme}{子群概形} $D'$，使 $D'_\eta=D'\otimes_{S'}\kappa(\eta)=H_\eta$。再由 EGA IV$_3$，9.7.7、9.3.3，可设 $D'$ 的\Term{fiber}{纤维}\Term{geometrically-reduced}{几何约化}。在（a）中，再次由 EGA IV$_3$，9.7.7、9.3.3，可设 $D'$ 的\Term{fiber}{纤维}\Term{connected}{连通}，因而\Term{geometrically-connected}{几何连通}（参见第 VI$_A$ 讲 2.1.1）。在（b）中，由 10.4，可设 $D'$ 是 $G'$ 的\Term{normal-subgroup}{正规子群}。

此外，在 7.8 的证明中已看到，存在整数 $n$，使 $\nu_\eta^n((A_\eta\times_{\kappa(\eta)}B_\eta)^n)=D'_\eta$，其中 $\nu_\eta,\nu_\eta^n$ 按 7.2.2（a）及 7.1（ii）定义。可用相同公式定义\Term{morphism}{态射}
\[
\nu':A'\times_{S'}B'\longrightarrow G',\qquad
\nu'{}^n:(A'\times_{S'}B')^n\longrightarrow G',
\]
且 $\nu'{}^n\otimes_{S'}\kappa(\eta)=\nu_\eta^n$。
\SourcePage{408}
因此由 10.1，可选择 $S'$，使\Term{morphism}{态射} $\nu'{}^n$ \Term{flat}{平坦}且经 $D'$ \Term{resolution}{分解}，并使这样得到的\Term{morphism}{态射} $\nu'{}^n:(A'\times_{S'}B')^n\to D'$ 为\Term{surjectivity}{满射}。于是由 7.5，\Term{morphism}{态射}\footnote{编者注：将下式中的 $n+1$ 改为 $2n$。}
\[
\nu'{}^{2n}=\mu\circ(\nu'{}^n\times_{S'}\nu'{}^n):(A'\times_{S'}B')^{2n}\longrightarrow D'
\]
相对于 $(\mathrm{fppf})$ \Term{topology}{拓扑}为\Term{covering}{覆盖}。故由 7.6，$D'$ 表示 $G'$ 中 $A',B'$ 的\Term{commutator}{交换子}\Term{presheaf}{预层}所伴的 $(\mathrm{fppf})$ \Term{sheaf}{层}。

此外，对每个 $s\in S'$，$\nu'{}^n$ 诱导\Term{surjective-morphism}{满射态射} $\nu_s^n:(A_s\times_{\kappa(s)}B_s)^n\to D'_s$。\footnote{编者注：以下对原文作了详细展开。}于是 $D'_s$ 是 $G_s$ 中包含 $\nu_s(A_s\times_{\kappa(s)}B_s)$ 的\Term{closed-subgroup}{闭子群}，故也包含 $(A_s,B_s)$。因 $\nu_s^n$ \Term{surjectivity}{满射}，$D'_s$ 等于 $(A_s,B_s)$，并由 7.6 表示 $G_s$ 中 $A_s,B_s$ 的\Term{commutator}{交换子} $(\mathrm{fppf})$ \Term{sheaf}{层}。
\end{proof}

\begin{SGAresult}{推论}{10.12.1}
\footnote{编者注：补入此推论；8.4 的证明要用到它。}设 $S$ 为\Term{integral}{整}\Term{scheme}{概形}，$\eta$ 为其\Term{generic-point}{泛点}，$G$ 为\Term{fiber}{纤维}\Term{smooth}{光滑}的\Term{finite-presentation}{有限表示} $S$-\Term{group-scheme}{群概形}。置 $K_\eta^0=G_\eta$，对每个 $p\in\mathbb N^*$ 置
\[
K_\eta^p=(K_\eta^{p-1},K_\eta^{p-1})
\quad\text{（相应地，}K_\eta^p=(G,K_\eta^{p-1})\text{）}.
\]
固定 $n\in\mathbb N^*$。则存在 $S$ 的\Term{nonempty}{非空}\Term{open-subset}{开集} $S'$ 及 $G|_{S'}$ 的\Term{normal-subgroup}{正规}\Term{subgroup-scheme}{子群概形} $D$，\Term{finite-presentation}{有限表示}且\Term{fiber}{纤维}\Term{smooth}{光滑}，使对每个 $s\in S'$ 都有 $D_s=K_s^n$。
\end{SGAresult}
对 $n$ 归纳，由 10.12 即得此结论。

\begin{SGAresult}{推论}{10.13}
设 $S$ 为\Term{integral}{整}\Term{scheme}{概形}，$\eta$ 为其\Term{generic-point}{泛点}，$G$ 为 $S$-\Term{group-scheme}{群概形}，$H$ 为 $G$ 的\Term{normal-subgroup}{正规} $S$-\Term{subgroup-scheme}{子群概形}；假设 $G,H$ 在 $S$ 上\Term{finite-presentation}{有限表示}，且\Term{generic-fibre}{泛纤维}\Term{smooth}{光滑}。\footnote{编者注：原文假设 $G,H$ \Term{finite-presentation}{有限表示}而 $H$ 的\Term{fiber}{纤维}\Term{smooth}{光滑}；为能应用 10.12，修改了假设，并展开了证明。}

若 $(H_\eta,H_\eta)=H_\eta$（相应地，$(G_\eta,H_\eta)=H_\eta$），则存在 $S$ 的\Term{nonempty}{非空}\Term{open-subset}{开集} $S'$，使对每个 $s\in S'$ 都有 $(H_s,H_s)=H_s$（相应地，$(G_s,H_s)=H_s$）。
\end{SGAresult}
\begin{proof}
事实上，由 10.12 的证明，存在 $S$ 的\Term{nonempty}{非空}\Term{open-subset}{开集} $S'$ 及 $G|_{S'}$ 的 $S'$-\Term{subgroup-scheme}{子群概形} $D$，\Term{finite-presentation}{有限表示}且\Term{fiber}{纤维}\Term{smooth}{光滑}，使对每个 $s\in S'$ 都有 $D_s=(H_s,H_s)$（相应地，$D_s=(G_s,H_s)$）。另一方面，因 $D_\eta=H_\eta$ 且 $D,H$ 在 $S'$ 上\Term{finite-presentation}{有限表示}，由 EGA IV$_3$，8.8.2.5，存在 $S'$ 的\Term{nonempty}{非空}\Term{open-subset}{开集} $S''$，使 $D|_{S''}=H|_{S''}$。故对每个 $s\in S''$，有 $H_s=D_s=(H_s,H_s)$（相应地，$H_s=D_s=(G_s,H_s)$）。
\end{proof}

\textbf{10.14。}关于\Term{finite-presentation}{有限表示}的 $S$-\Term{group-scheme}{群概形}\Term{category}{范畴}的陈述 10.2、10.3，可推广到如下二元组的\Term{category}{范畴}：一个\Term{finite-presentation}{有限表示}的 $S$-\Term{group-scheme}{群概形} $G$，加上一个以 $G$ 为\Term{operator-group}{算子群}的\Term{finite-presentation}{有限表示} $S$-\Term{scheme}{概形}。准确地说，在 10.0 开头所回顾的情形中：

（i）设 $j\in I$，$G_j,G'_j$ 是两个\Term{finite-presentation}{有限表示}的 $S_j$-\Term{group-scheme}{群概形}，$H_j$（相应地，$H'_j$）是以 $G_j$（相应地，$G'_j$）为\Term{operator-group}{算子群}的\Term{finite-presentation}{有限表示} $S_j$-\Term{scheme}{概形}。对 $i\in I$、$i\geq j$，置 $G_i=G_j\times_{S_j}S_i$、$G=G_j\times_{S_j}S$，并类似地定义 $G'_i,G',H_i,H,H'_i,H'$。记 $\operatorname{Dihom}_{S\text{-gr.}}((G,H),(G',H'))$ 为从二元组 $(G,H)$ 到二元组 $(G',H')$ 的 $S$-群与带\Term{operator-group}{算子群}的 $S$-\Term{scheme}{概形}的\Term{dimorphism}{双态射}所成的\Term{set}{集合}。
\SourcePage{409}
则\Term{canonical}{典范}映射
\[
\varinjlim_{i\geq j}\operatorname{Dihom}_{S_i\text{-gr.}}((G_i,H_i),(G'_i,H'_i))
\longrightarrow\operatorname{Dihom}_{S\text{-gr.}}((G,H),(G',H'))
\]
是\Term{bijection}{双射}。

（ii）设 $G$ 是\Term{finite-presentation}{有限表示}的 $S$-\Term{group-scheme}{群概形}，$H$ 是以 $G$ 为\Term{operator-group}{算子群}的\Term{finite-presentation}{有限表示} $S$-\Term{scheme}{概形}。则存在指标 $j\in I$、\Term{finite-presentation}{有限表示}的 $S_j$-\Term{group-scheme}{群概形} $G_j$、以 $G_j$ 为\Term{operator-group}{算子群}的\Term{finite-presentation}{有限表示} $S_j$-\Term{scheme}{概形} $H_j$，以及从 $(G_j\times_{S_j}S,H_j\times_{S_j}S)$ 到 $(G,X)$ 的 $S$-群与带\Term{operator-group}{算子群}的 $S$-\Term{scheme}{概形}的\Term{diisomorphism}{双同构}。

\begin{SGAresult}{定义}{10.15}
\footnote{编者注：通过区分“\Term{formally-homogeneous}{形式齐性}”与“\Term{homogeneous}{齐性}”\Term{object}{对象}的概念修正了原文；参见第 IV 讲 §§5.1、6.7。}设 $T$ 为 $(\mathrm{Sch}/S)$ 上不细于\Term{canonical-topology}{典范拓扑}的一个\Term{topology}{拓扑}。给定 $S$-\Term{group-scheme}{群概形} $G$ 及以 $G$ 为\Term{operator-group}{算子群}的 $S$-\Term{scheme}{概形} $X$，若\Term{morphism}{态射}
\[
\Phi:G\times_S X\longrightarrow X\times_S X,
\qquad(g,x)\longmapsto(gx,x)
\]
（对每个 $S'\to S$ 及 $g\in G(S'),x\in X(S')$ 按此式定义）在\Term{topology}{拓扑} $T$ 的\Term{sheaf}{层}\Term{category}{范畴}中是\Term{epimorphism}{满态射}，就称 $X$ 为 $G$ 下的\Term{formally-homogeneous}{形式齐性空间}（相对于\Term{topology}{拓扑} $T$）。这\Term{equivalence}{等价}于说 $\Phi$ 相对于\Term{topology}{拓扑} $T$ 为\Term{covering}{覆盖}（参见第 IV 讲 4.4.3）。

若 $X$ 为\Term{formally-homogeneous}{形式齐性}的，且\Term{morphism}{态射} $X\to S$ 也相对于\Term{topology}{拓扑} $T$ 为\Term{covering}{覆盖}，就称 $X$ 为\Term{homogeneous-space}{齐性空间}。

特别地，若 $\Phi$ 为\Term{isomorphism}{同构}，就称 $X$ 为 $G$ 下的\Term{formally-principal-homogeneous}{形式主齐性空间}；若 $\Phi$ 为\Term{isomorphism}{同构}且\Term{morphism}{态射} $X\to S$ 相对于\Term{topology}{拓扑} $T$ 为\Term{covering}{覆盖}，就称 $X$ 为\Term{principal-homogeneous-bundle}{主齐性丛}（或 $G$-\Term{torsor}{挠子}）（参见第 IV 讲 5.1.5 及 5.1.6（ii））。
\end{SGAresult}

\begin{SGAresult}{命题}{10.16}
处在 10.0 开头考察的情形中。设 $j\in I$，$G_j$ 为 $S_j$-\Term{group-scheme}{群概形}，$X_j$ 为以 $G_j$ 为\Term{operator-group}{算子群}的 $S_j$-\Term{scheme}{概形}。假设 $G_j,X_j$ 在 $S_j$ 上\Term{finite-presentation}{有限表示}。

则 $X=X_j\times_{S_j}S$ 是 $G=G_j\times_{S_j}S$ 下相对于 $(\mathrm{fppf})$ \Term{topology}{拓扑}的\Term{homogeneous-space}{齐性空间}（相应地，\Term{principal-homogeneous-bundle}{主齐性丛}）的充要条件是，存在指标 $i\geq j$，使 $X_i=X_j\times_{S_j}S_i$ 是 $G_i=G_j\times_{S_j}S_i$ 下的\Term{homogeneous-space}{齐性空间}（相应地，\Term{principal-homogeneous-bundle}{主齐性丛}）。
\end{SGAresult}
结合 10.14 及 EGA IV$_3$，8.8.2、8.8.3、8.10.5，此陈述由 10.1 中回顾的关于 $(\mathrm{fppf})$ \Term{topology}{拓扑}的\Term{covering}{覆盖}\Term{morphism}{态射}的性质\Term{pushout}{推出}。\footnote{编者注：关于用 \Term{cech-cohomology}{Čech 上同调}表述的\Term{torsor}{挠子}的\Term{limit}{极限}性质，另见 SGA 4，第 VII 讲 5.7 及其推论。文章 [Ma07] 对此作了详细说明。}

\SGAsection{11}{仿射群概形}
\SourcePage{410}
\textbf{11.0。回顾。}\footnote{编者注：补入了这一回顾段。}设 $q:X\to S$ 是\Term{quasi-compact}{拟紧}且\Term{quasi-separated}{拟分离}的\Term{scheme}{概形}\Term{morphism}{态射}（参见 EGA IV$_1$，1.1、1.2），$\mathcal F$ 是\Term{quasi-coherent}{拟凝聚} $\mathcal O_X$-模。回顾，$q_*(\mathcal F)$ 为\Term{quasi-coherent}{拟凝聚} $\mathcal O_S$-模（EGA I，9.2.1）。此外，由 EGA III，1.4.15（加上 EGA IV$_1$，1.7.21 的补充），有下述（c），而该处的证明也给出（a）、（b）：

（a）若 $\mathcal F$ 是\Term{quasi-coherent}{拟凝聚}\Term{submodule}{子模} $\mathcal F_\alpha$ 的\Term{filtered}{滤过}\Term{inductive-limit}{归纳极限}，则 $q_*(\mathcal F)=\varinjlim_\alpha q_*(\mathcal F_\alpha)$。

（b）若 $\mathcal E$ 是\Term{flat}{平坦} $\mathcal O_S$-模，则\Term{canonical-morphism}{典范态射} $\mathcal E\otimes_{\mathcal O_S}q_*(\mathcal O_X)\to q_*q^*(\mathcal E)$ 为\Term{isomorphism}{同构}。

（c）设 $p:S'\to S$ 是\Term{flat}{平坦}\Term{morphism}{态射}，$q':X'\to S'$ 由 $q$ 经\Term{base-change}{基变换}得到，$\mathcal F'$ 为 $\mathcal F$ 在 $X'$ 上的\Term{inverse-image}{逆像}。则\Term{canonical-morphism}{典范态射} $p^*q_*(\mathcal F)\to q'_*(\mathcal F')$ 为\Term{isomorphism}{同构}。

事实上，设 $U=\operatorname{Spec}(A)$ 是 $S$ 的任意\Term{affine}{仿射}\Term{open-subset}{开集}。由假设，$q^{-1}(U)$ 是\Term{affine}{仿射}\Term{open-subset}{开集} $V_i=\operatorname{Spec}(B_i)$（$i=1,\ldots,n$）之并，且每个交集 $V_i\cap V_j$ 都是\Term{finite}{有限}多个\Term{affine}{仿射}\Term{open-subset}{开集} $W_{ijk}=\operatorname{Spec}(C_{ijk})$ 之并。于是 $\Gamma(U,q_*(\mathcal F))=\Gamma(q^{-1}(U),\mathcal F)$ 是\Term{morphism}{态射}
\[
\bigoplus_{i=1}^n\Gamma(V_i,\mathcal F)\longrightarrow\bigoplus_{i,j,k}\Gamma(W_{ijk},\mathcal F)
\]
的核。由此得（a），因为上式各项都与\Term{filtered}{滤过}\Term{inductive-limit}{归纳极限}交换（$V_i,W_{ijk}$ 为\Term{affine}{仿射}的，故\Term{quasi-compact}{拟紧}）。证明（b）：$E=\Gamma(U,\mathcal E)$ 是\Term{flat}{平坦} $A$-模，而 $\Gamma(U,q_*q^*(\mathcal E))$ 是\Term{morphism}{态射}
\[
\bigoplus_{i=1}^n B_i\otimes_A E\longrightarrow\bigoplus_{i,j,k}C_{ijk}\otimes_A E
\]
的核 $K(E)$。因 $E$ 在 $A$ 上\Term{flat}{平坦}，此核等同于 $K(A)\otimes_A E=\mathcal O_X(q^{-1}(U))\otimes_A E$。最后，若 $U'$ 为 $S'$ 中位于 $U$ 上方的任意\Term{affine}{仿射}\Term{open-subset}{开集}，则 $A'=\mathcal O(U')$ 为\Term{flat}{平坦} $A$-\Term{algebra}{代数}，且如上可得
\[
\mathcal F(q^{-1}(U))\otimes_A A'\xrightarrow{\sim}\mathcal F'(q'{}^{-1}(U')).
\]

\textbf{记号。}设 $S$ 为\Term{scheme}{概形}，$X$ 为 $S$-\Term{scheme}{概形}，$f:X\to S$ 为\Term{structural-morphism}{结构态射}；记 $\mathcal A(X)=f_*(\mathcal O_X)$。
\begin{SGAresult}{\Term{lemma}{引理}}{11.1}
设 $X,Y$ 是 $S$ 上\Term{quasi-compact}{拟紧}且\Term{quasi-separated}{拟分离}的两个 $S$-\Term{scheme}{概形}，$f:X\to S$、$g:Y\to S$ 是\Term{structural-morphism}{结构态射}。则\Term{canonical-homomorphism}{典范同态}
\[
\varphi:\mathcal A(X)\otimes_{\mathcal O_S}\mathcal A(Y)\longrightarrow\mathcal A(X\times_S Y)
\]
在下列各情形中为\Term{isomorphism}{同构}：\footnote{编者注：注意，若 $k$ 为域，而 $X$ 为无穷多个 $S=\operatorname{Spec}k$ 的副本之和（从而 $X$ 不\Term{quasi-compact}{拟紧}），则 $\mathcal A(X)=k^X$，\Term{canonical-morphism}{典范态射} $k^X\otimes k^X\to k^{X\times X}$ 并非\Term{surjectivity}{满射}。}

 a）$f,g$ 为\Term{affine}{仿射}\Term{morphism}{态射}；

 b）$g$（或 $f$）\Term{flat}{平坦}且\Term{affine}{仿射}；
\SourcePage{411}

 c）$g$ \Term{flat}{平坦}，且 $f_*(\mathcal O_X)$ 是\Term{flat}{平坦} $\mathcal O_S$-模。
\end{SGAresult}
\begin{proof}
在（b）中，假设 $g$ \Term{flat}{平坦}且\Term{affine}{仿射}。置 $S'=\operatorname{Spec}\mathcal A(X)$，$Y'=Y\times_S S'$，$g'=g\times_S S'$，并记 $v$ 为\Term{morphism}{态射} $S'\to S$：
\[
\begin{tikzcd}
Y'\arrow[r]\arrow[d,"g'"']&Y\arrow[d,"g"]\\
\operatorname{Spec}\mathcal A(X)=S'\arrow[r,"v"']&S.
\end{tikzcd}
\]
在（a）、（b）中，$g$ 为\Term{affine}{仿射}\Term{morphism}{态射}，故由 EGA II，1.5.2 有
\[
g'_*(\mathcal O_{Y'})=v^*g_*(\mathcal O_Y)=\mathcal A(Y)\otimes_{\mathcal O_S}\mathcal O_{S'}.\tag{1}
\]
在（c）中，由 11.0（c）亦有此等式，因为 $S'$ 在 $S$ 上\Term{flat}{平坦}，而 $g$ \Term{quasi-compact}{拟紧}且\Term{quasi-separated}{拟分离}。

另一方面（EGA II，1.2.7），$f:X\to S$ 借助\Term{morphism}{态射} $p:X\to S'$ 经 $v$ \Term{resolution}{分解}，并有 $p_*(\mathcal O_X)=\mathcal O_{S'}$ 及 $X\times_S Y=X\times_{S'}Y'$。因 $f$ \Term{quasi-separated}{拟分离}，$p$ 亦然（EGA IV$_1$，1.2.2）；又因 $f$ \Term{quasi-compact}{拟紧}而 $v$ \Term{quasi-separated}{拟分离}，$p$ 也\Term{quasi-compact}{拟紧}（EGA IV$_1$，1.2.4）。考虑\Term{cartesian}{笛卡尔}方块
\[
\begin{tikzcd}
X\times_S Y\arrow[r,"p'"]\arrow[d]&Y'\arrow[d,"g'"]\\
X\arrow[r,"p"']&S'.
\end{tikzcd}
\]
在（b）、（c）中，$Y$ 在 $S$ 上\Term{flat}{平坦}，故 $Y'$ 在 $S'$ 上\Term{flat}{平坦}；再次应用 11.0（c），得
\[
p'_*(\mathcal O_{X\times_S Y})=g'{}^*p_*(\mathcal O_X)=g'{}^*(\mathcal O_{S'})=\mathcal O_{Y'}.\tag{2}
\]
在（a）中亦有此等式，因为此时 $p,p'$ 都是\Term{isomorphism}{同构}。

最后，因 $v$ 是\Term{affine}{仿射}\Term{morphism}{态射}，由 EGA II，1.4.7，对每个\Term{quasi-coherent}{拟凝聚} $\mathcal O_S$-模 $\mathcal F$，有 $v_*(\mathcal F\otimes_{\mathcal O_S}\mathcal O_{S'})=\mathcal F\otimes_{\mathcal O_S}\mathcal A(X)$。结合（2）、（1），得
\[
\begin{aligned}
\mathcal A(X\times_S Y)
&=v_*g'_*p'_*(\mathcal O_{X\times_S Y})
\overset{(2)}=v_*g'_*(\mathcal O_{Y'})\\
&\overset{(1)}=v_*(\mathcal A(Y)\otimes_{\mathcal O_S}\mathcal O_{S'})
=\mathcal A(Y)\otimes_{\mathcal O_S}\mathcal A(X).
\end{aligned}
\]
\end{proof}
\begin{SGAresult}{推论}{11.2}
\Term{functor}{函子} $X\mapsto\operatorname{Spec}\mathcal A(X)$，从 $(\mathrm{Sch}/S)$ 中由满足 $X$ 在 $S$ 上\Term{flat}{平坦}、\Term{quasi-compact}{拟紧}、\Term{quasi-separated}{拟分离}，且 $\mathcal A(X)$ 是\Term{flat}{平坦} $\mathcal O_S$-模的 $S$-\Term{scheme}{概形}构成的\Term{full-subcategory}{满子范畴}，到 $S$ 上\Term{flat}{平坦}且\Term{affine}{仿射}的 $S$-\Term{scheme}{概形}构成的\Term{full-subcategory}{满子范畴}，与\Term{finite-product}{有限积}交换，因而将 $S$-\Term{group-scheme}{群概形}映为 $S$-\Term{group-scheme}{群概形}。
\end{SGAresult}
\begin{SGAresult}{定义}{11.3}
给定在 $S$ 上\Term{flat}{平坦}、\Term{quasi-compact}{拟紧}、\Term{quasi-separated}{拟分离}，且 $\mathcal A(G)$ 在 $\mathcal O_S$ 上\Term{flat}{平坦}的 $S$-\Term{group-scheme}{群概形} $G$，\footnote{编者注：若 $S$ 为\Term{dimension}{维数} $\leq2$ 的\Term{regular}{正则}\Term{locally-noetherian}{局部诺特}\Term{scheme}{概形}，而 $X$ 是在 $S$ 上\Term{flat}{平坦}、\Term{quasi-compact}{拟紧}、\Term{quasi-separated}{拟分离}的 $S$-\Term{scheme}{概形}，则 $\mathcal A(X)$ 是\Term{flat}{平坦} $\mathcal O_S$-模；参见 [Ray70a]，VII 3.2。}将 $S$-\Term{group-scheme}{群概形} $G^{\mathrm{af}}=\operatorname{Spec}\mathcal A(G)$ 称为 $G$ 的\Term{affine-envelope}{仿射包络}，记为 $G^{\mathrm{af}}$。
\end{SGAresult}

\SourcePage{412}
\begin{SGAresult}{命题}{11.3.1}
\footnote{编者注：补入了本命题；关于\Term{morphism}{态射} $G\to G^{\mathrm{af}}$ 及其核的研究，另见后面补入的第 12 节。}\Term{canonical-morphism}{典范态射} $\tau_G:G\to G^{\mathrm{af}}$ 是 $S$-\Term{group-morphism}{群态射}，且具有如下\Term{universal-property}{泛性质}：

（i）对每个 $S$-\Term{scheme}{概形}\Term{morphism}{态射} $\varphi:G\to H$，其中 $H$ 在 $S$ 上\Term{affine}{仿射}，存在唯一的 $S$-\Term{scheme}{概形}\Term{morphism}{态射} $\varphi':G^{\mathrm{af}}\to H$，使 $\varphi=\varphi'\circ\tau_G$。

（ii）若 $H$ 还是 $S$-\Term{group-scheme}{群概形}，且 $\varphi$ 为 $S$-\Term{group-morphism}{群态射}，则 $\varphi'$ 亦然。
\end{SGAresult}

\textbf{11.4。}\footnote{编者注：11.4--11.6 中使用了 $\operatorname{Hom}_{\mathcal O_S}(\mathbf W(\mathcal E),\mathbf W(\mathcal F))$，而非 $\operatorname{Hom}_{\mathcal O_S}(\mathbf V(\mathcal F),\mathbf V(\mathcal E))$（参见第 I 讲 4.6.3），并利用补入的 11.0（b）简化了原文。}设 $\mathcal E,\mathcal F$ 为两个\Term{quasi-coherent}{拟凝聚} $\mathcal O_S$-模。考虑 $S$-\Term{functor}{函子} $\operatorname{Hom}_{\mathcal O_S}(\mathbf W(\mathcal E),\mathbf W(\mathcal F))$（参见第 I 讲 3.1.4），即对每个 $S$-\Term{scheme}{概形} $f:X\to S$，
\[
\operatorname{Hom}_{\mathcal O_S}(\mathbf W(\mathcal E),\mathbf W(\mathcal F))(X)
=\operatorname{Hom}_{\mathcal O_X}(\mathbf W(\mathcal E)_X,\mathbf W(\mathcal F)_X).
\]
此外，由第 I 讲 4.6.2，$\mathbf W(\mathcal E)_X=\mathbf W(f^*(\mathcal E))$（对 $\mathcal F$ 亦然），并且
\[
\operatorname{Hom}_{\mathcal O_X}(\mathbf W(f^*(\mathcal E)),\mathbf W(f^*(\mathcal F)))
=\operatorname{Hom}_{\mathcal O_X}(f^*(\mathcal E),f^*(\mathcal F)).
\]
因此得到（最后一个等式使用\Term{adjunction}{伴随}公式）
\[
\begin{aligned}
\operatorname{Hom}_{\mathcal O_S}(\mathbf W(\mathcal E),\mathbf W(\mathcal F))(X)
&=\operatorname{Hom}_{\mathcal O_X}(f^*(\mathcal E),f^*(\mathcal F))\\
&=\operatorname{Hom}_{\mathcal O_S}(\mathcal E,f_*f^*(\mathcal F)).
\end{aligned}\tag{$\dagger$}
\]
\begin{SGAresult}{命题}{11.5}
设 $X$ 是 $S$ 上\Term{quasi-compact}{拟紧}且\Term{quasi-separated}{拟分离}的 $S$-\Term{scheme}{概形}，$f:X\to S$ 为\Term{structural-morphism}{结构态射}，$\mathcal E,\mathcal E'$ 为两个\Term{quasi-coherent}{拟凝聚} $\mathcal O_S$-模。假设满足下列两个条件之一：

 a）$f$ 为\Term{affine}{仿射}\Term{morphism}{态射}；

 b）$\mathcal E$ 在 $\mathcal O_S$ 上\Term{flat}{平坦}。

则\Term{canonical-morphism}{典范态射} $\mathcal E\otimes_{\mathcal O_S}\mathcal A(X)\to f_*f^*(\mathcal E)$ 为\Term{isomorphism}{同构}，因此有
\[
\operatorname{Hom}_{\mathcal O_S}(\mathbf W(\mathcal E'),\mathbf W(\mathcal E))(X)
=\operatorname{Hom}_{\mathcal O_S}(\mathcal E',\mathcal E\otimes_{\mathcal O_S}\mathcal A(X)).
\]
\end{SGAresult}
事实上，第二个断言由 11.4 及第一个断言得到；第一个断言在（a）中由 EGA II，1.4.7 得到，在（b）中由 11.0（b）得到。\footnote{编者注：简化了原文。原文使用\Term{isomorphism}{同构} $\operatorname{Hom}_{\mathcal O_S}(\mathbf W(\mathcal E'),\mathbf W(\mathcal E))\simeq\operatorname{Hom}_{\mathcal O_S}(\mathbf V(\mathcal E),\mathbf V(\mathcal E'))$，继而将右边\Term{embedding}{嵌入}
$\operatorname{Hom}_S(\mathbf V(\mathcal E),\mathbf V(\mathcal E'))=\operatorname{Hom}_{\mathcal O_S}(\mathcal E',f_*f^*(\operatorname{Sym}(\mathcal E)))$，并对 $\mathbf V(\mathcal E)=\operatorname{Spec}(\operatorname{Sym}(\mathcal E))$ 应用 EGA III，4.1.15，以\Term{pushout}{推出} 11.0（b）。}

\textbf{11.6。}\footnote{编者注：展开了 11.6，并将所得结果列为命题 11.6.1。}设 $G$ 为 $S$-\Term{group-scheme}{群概形}，$\mathcal E$ 为\Term{quasi-coherent}{拟凝聚} $\mathcal O_S$-模。在 $\mathcal E$ 上给定 $G$-$\mathcal O_S$-模结构（即 $G$ 在 $\mathbf W(\mathcal E)$ 上的 $\mathcal O_S$-\Term{linear-action}{线性作用}，参见第 I 讲 4.7.1），\Term{equivalence}{等价}于给定 $S$-\Term{monoid-functor}{幺半群函子}的\Term{morphism}{态射}
\[
\rho:G\longrightarrow\operatorname{End}_{\mathcal O_S}(\mathbf W(\mathcal E))
\]
（事实上，这样的 $\rho$ 必将 $G$ 映入 $\operatorname{Aut}_{\mathcal O_S}(\mathbf W(\mathcal E))$）。

\SourcePage{413}
由 11.4，给定 $S$-\Term{morphism-of-functors}{函子态射} $\rho:G\to\operatorname{End}_{\mathcal O_S}(\mathbf W(\mathcal E))$，\Term{equivalence}{等价}于给定 $\operatorname{Hom}_{\mathcal O_G}(f^*(\mathcal E),f^*(\mathcal E))$ 的元素 $\theta$，它通过\Term{adjunction}{伴随}对应于 $\operatorname{Hom}_{\mathcal O_S}(\mathcal E,f_*f^*(\mathcal E))$ 的元素 $\delta$，其中 $f$ 表示\Term{projection}{投影} $G\to S$。

设 $m:G\times_S G\to G$ 为\Term{multiplication}{乘法}，$\delta_G$ 为\Term{morphism}{态射} $\mathcal O_G\to m_*(\mathcal O_{G\times_S G})$，$\phi$ 为\Term{projection}{投影} $G\times_S G\to S$（等于 $f\circ m$）。用 $\boxtimes$ 表示“外”\Term{tensor-product}{张量积}很方便；这样得到\Term{morphism}{态射}
\[
\operatorname{id}_{\mathcal E}\boxtimes\delta_G:
 f^*(\mathcal E)=\mathcal E\boxtimes_{\mathcal O_S}\mathcal O_G
\longrightarrow\mathcal E\boxtimes_{\mathcal O_S}m_*(\mathcal O_{G\times G}).
\]
稍滥用记号，仍用 $\operatorname{id}_{\mathcal E}\boxtimes\delta_G$ 表示前一\Term{morphism}{态射}与\Term{canonical-morphism}{典范态射}
\[
\mathcal E\boxtimes_{\mathcal O_S}m_*(\mathcal O_{G\times G})
\longrightarrow m_*m^*f^*(\mathcal E)=m_*\phi^*(\mathcal E)
\]
的复合。

另一方面，用 $h:G\to S$ 表示 $f:G\to S$ 的另一副本，考虑如下\Term{commutative-diagram}{交换图}，其中 $p,q$ 为两个\Term{projection}{投影}：
\[
\begin{tikzcd}
G\times_S G\arrow[r,"q"]\arrow[d,"p"']\arrow[dr,"\phi" description]&G\arrow[d,"h"]\\
G\arrow[r,"f"']&S.
\end{tikzcd}
\]
仍记 $\delta$ 为\Term{morphism}{态射} $\mathcal E\to h_*h^*(\mathcal E)$，得到\Term{morphism}{态射}
\[
\begin{aligned}
\delta\boxtimes\operatorname{id}_{\mathcal O_G}:
 f^*(\mathcal E)&=\mathcal E\boxtimes_{\mathcal O_S}\mathcal O_G\\
&\longrightarrow h_*h^*(\mathcal E)\boxtimes_{\mathcal O_S}\mathcal O_G
=f^*h_*h^*(\mathcal E),
\end{aligned}
\]
并仍滥用记号，用 $\delta\boxtimes\operatorname{id}_{\mathcal O_G}$ 表示此\Term{morphism}{态射}与\Term{canonical-morphism}{典范态射} $f^*h_*h^*(\mathcal E)\to p_*\phi^*(\mathcal E)$ 的复合。

那么，$\rho$ 与\Term{multiplication}{乘法}相容的条件\Term{equivalence}{等价}于说，对 $S$ 的每个\Term{open-subset}{开集} $U$，下图交换：
\[
\begin{tikzcd}[column sep=small]
\Gamma(U,\mathcal E)\arrow[r,"\delta"]\arrow[d,"\delta"']&
\Gamma(U\times_S G,\mathcal E\boxtimes_{\mathcal O_S}\mathcal O_G)
\arrow[d,"\operatorname{id}_{\mathcal E}\boxtimes\delta_G"]\\
\Gamma(U\times_S G,\mathcal E\boxtimes_{\mathcal O_S}\mathcal O_G)
\arrow[r,"\delta\boxtimes\operatorname{id}_{\mathcal O_G}"']&
\Gamma(U\times_S G\times_S G,\mathcal E\boxtimes_{\mathcal O_S}\mathcal O_G\boxtimes_{\mathcal O_S}\mathcal O_G).
\end{tikzcd}\tag{1}
\]
此外，\Term{unit-section}{单位截面} $\varepsilon:S\to G$ 诱\Term{derived-morphism}{导态射} $u:\mathcal O_G\to\varepsilon_*\varepsilon^*(\mathcal O_G)=\varepsilon_*(\mathcal O_S)$，而 $\rho$ 保持\Term{identity-element}{单位元}的条件\Term{equivalence}{等价}于下图交换：
\[
\begin{tikzcd}[column sep=small]
\Gamma(U,\mathcal E)\arrow[rr,"\delta"]\arrow[dr,"\sim"']&&
\Gamma(U\times_S G,\mathcal E\boxtimes_{\mathcal O_S}\mathcal O_G)
\arrow[dl,"\operatorname{id}_{\mathcal E}\boxtimes u"]\\
&\Gamma(U\times_S S,\mathcal E\boxtimes_{\mathcal O_S}\mathcal O_S).&
\end{tikzcd}\tag{2}
\]
由此可见，在 $\mathcal E$ 上给定 $G$-$\mathcal O_S$-模结构，\Term{equivalence}{等价}于给定满足上述条件（1）、（2）的 $\mathcal O_S$-\Term{module-morphism}{模态射} $\delta:\mathcal E\to f_*f^*(\mathcal E)$。此时，由 $\delta$ 经\Term{adjunction}{伴随}得到的\Term{morphism}{态射} $\theta:f^*(\mathcal E)\to f^*(\mathcal E)$ 是\Term{isomorphism}{同构}，
\SourcePage{414}
因为它对应于\Term{set}{集合}意义下由 $(g,x)\mapsto(g,gx)$ 定义的\Term{isomorphism}{同构} $G\times_S\mathbf W(\mathcal E)\xrightarrow{\sim}G\times_S\mathbf W(\mathcal E)$；另见第 I 讲 6.5.4。

现在设 $G$ 在 $S$ 上\Term{flat}{平坦}、\Term{quasi-compact}{拟紧}、\Term{quasi-separated}{拟分离}，且 $\mathcal A(G)$ 是\Term{flat}{平坦} $\mathcal O_S$-模。由 11.1（c），\Term{canonical-morphism}{典范态射} $\mathcal A(G)\otimes_{\mathcal O_S}\mathcal A(G)\to\mathcal A(G\times_S G)$ 为\Term{isomorphism}{同构}；将\Term{morphism}{态射} $\delta_G:\mathcal A(G)\to\mathcal A(G)\otimes_{\mathcal O_S}\mathcal A(G)$ 记为 $\Delta$。

若还有 $\mathcal E\otimes_{\mathcal O_S}\mathcal A(G)=f_*f^*(\mathcal E)$（由 11.5，当 $G\to S$ 为\Term{affine}{仿射}\Term{morphism}{态射}，或 $\mathcal E$ 在 $\mathcal O_S$ 上\Term{flat}{平坦}时即如此），则条件（1）、（2）\Term{equivalence}{等价}于下列条件，它们表明 $\delta:\mathcal E\to\mathcal E\otimes_{\mathcal O_S}\mathcal A(G)$ 赋予 $\mathcal E$ 一个右 $\mathcal A(G)$-\Term{comodule}{余模}结构（参见第 I 讲 4.7.2）：

（CM 1）置 $\mathcal A=\mathcal A(G)$，下图交换：
\[
\begin{tikzcd}
\mathcal E\arrow[r,"\delta"]\arrow[d,"\delta"']&\mathcal E\otimes\mathcal A\arrow[d,"\operatorname{id}_{\mathcal E}\otimes\Delta"]\\
\mathcal E\otimes\mathcal A\arrow[r,"\delta\otimes\operatorname{id}_{\mathcal A}"']&\mathcal E\otimes\mathcal A\otimes\mathcal A.
\end{tikzcd}
\]
（CM 2）记 $\eta:\mathcal A\to\mathcal O_S$ 为\Term{morphism}{态射} $\mathcal A(\varepsilon)$，则下图交换：
\[
\begin{tikzcd}
\mathcal E\arrow[rr,"\delta"]\arrow[dr,"\sim"']&&\mathcal E\otimes\mathcal A\arrow[dl,"\operatorname{id}_{\mathcal E}\otimes\eta"]\\
&\mathcal E\otimes\mathcal O_S.&
\end{tikzcd}
\]

\textbf{注记 11.6.A。}回顾，以 $\mathbf V(\mathcal E)$ 表示 $S$ 上表示\Term{functor}{函子} $\mathbf V(\mathcal E)$ 的\Term{vector-fibration}{向量纤维化}，即对每个 $S'\to S$，有 $\mathbf V(\mathcal E)(S')=\operatorname{Hom}_{\mathcal O_{S'}}(\mathcal E_{S'},\mathcal O_{S'})$。由第 I 讲 4.6.2，有 $S$-\Term{monoid-functor}{幺半群函子}的\Term{anti-isomorphism}{反同构}
\[
\operatorname{End}_{\mathcal O_S}(\mathbf W(\mathcal E))\simeq\operatorname{End}_{\mathcal O_S}(\mathbf V(\mathcal E)).
\]
故若 $\mathcal E$ 是左 $G$-$\mathcal O_S$-模，就有 $G$ 在 $\mathbf V(\mathcal E)$ 上的\Term{right-action}{右作用} $\mu:\mathbf V(\mathcal E)\times_S G\to\mathbf V(\mathcal E)$，它在\Term{set}{集合}意义下由 $(\varphi g)(x)=\varphi(gx)$ 定义，其中 $g\in G(S')$、$x\in\Gamma(S',\mathcal E_{S'})$、$\varphi\in\operatorname{Hom}_{\mathcal O_{S'}}(\mathcal E_{S'},\mathcal O_{S'})$。于是得到\Term{commutative-diagram}{交换图}
\[
\begin{tikzcd}[column sep=small]
\mathbf V(\mathcal E)\times_S G\times_S G\arrow[r,"\mu\times\operatorname{id}_G"]\arrow[d,"\operatorname{id}_{\mathbf V(\mathcal E)}\times m"']&\mathbf V(\mathcal E)\times_S G\arrow[d,"\mu"]\\
\mathbf V(\mathcal E)\times_S G\arrow[r,"\mu"']&\mathbf V(\mathcal E),
\end{tikzcd}
\qquad
\begin{tikzcd}[column sep=small]
\mathbf V(\mathcal E)&\mathbf V(\mathcal E)\times_S G\arrow[l,"\mu"']\\
&\mathbf V(\mathcal E)\times_S S\arrow[u,"\operatorname{id}_{\mathbf V(\mathcal E)}\times\varepsilon"']\arrow[ul,"\sim"].
\end{tikzcd}
\]
当 $G$ 在 $S$ 上\Term{flat}{平坦}、\Term{quasi-compact}{拟紧}、\Term{quasi-separated}{拟分离}，$\mathcal A(G)$ 是\Term{flat}{平坦} $\mathcal O_S$-模，且满足 11.5 的条件之一时，同样得到条件（CM 1）、（CM 2）。

因此，已经得到：
\begin{SGAresult}{命题}{11.6.1}
设 $G$ 是在 $S$ 上\Term{flat}{平坦}、\Term{quasi-compact}{拟紧}、\Term{quasi-separated}{拟分离}的 $S$-\Term{group-scheme}{群概形}，使 $\mathcal A(G)$ 为\Term{flat}{平坦} $\mathcal O_S$-模；设 $\mathcal E$ 为\Term{quasi-coherent}{拟凝聚} $\mathcal O_S$-模。
\SourcePage{415}

（i）给定 $\mathcal A(G)$-\Term{comodule}{余模}结构 $\delta:\mathcal E\to\mathcal E\otimes_{\mathcal O_S}\mathcal A(G)$，\Term{equivalence}{等价}于在 $\mathcal E$ 上给定 $G^{\mathrm{af}}$-$\mathcal O_S$-模结构（即 $G^{\mathrm{af}}$ 在 $\mathcal E$ 上的 $\mathcal O_S$-\Term{linear-action}{线性作用}）。与 $S$-\Term{group-morphism}{群态射} $G\to G^{\mathrm{af}}$ 复合，就在 $\mathcal E$ 上定义了 $G$-$\mathcal O_S$-模结构。

（ii）若 $\mathcal E$ 还\Term{flat}{平坦}，则 $G$ 在 $\mathcal E$ 上的每个 $\mathcal O_S$-\Term{linear-action}{线性作用}都经 $G^{\mathrm{af}}$ \Term{resolution}{分解}，并对应于 $\mathcal E$ 上唯一的 $\mathcal A(G)$-\Term{comodule}{余模}结构。
\end{SGAresult}

\begin{SGAresult}{\Term{lemma}{引理}}{11.7}
设 $G$ 是在 $S$ 上\Term{flat}{平坦}、\Term{quasi-compact}{拟紧}、\Term{quasi-separated}{拟分离}的 $S$-\Term{group-scheme}{群概形}，使 $\mathcal A=\mathcal A(G)$ 为\Term{flat}{平坦} $\mathcal O_S$-模。设 $\mathcal E$ 为\Term{quasi-coherent}{拟凝聚} $\mathcal O_S$-模，$\delta:\mathcal E\to\mathcal E\otimes\mathcal A$ 为 $\mathcal A$-\Term{comodule}{余模}结构，$\rho:G\to\operatorname{Aut}_{\mathcal O_S}(\mathbf W(\mathcal E))$ 为相应的 $G$ 在 $\mathcal E$ 上的\Term{action}{作用}。

设 $\mathcal E_0$ 为 $\mathcal E$ 的\Term{quasi-coherent}{拟凝聚} $\mathcal O_S$-\Term{submodule}{子模}，使 $\mu$ 在 $\mathcal E_0$ 上的\Term{restriction}{限制} $\delta_0$ 经 $\mathcal E_0\to\mathcal E_0\otimes\mathcal A$ \Term{resolution}{分解}，即有\Term{commutative-diagram}{交换图}
\[
\begin{tikzcd}
\mathcal E_0\arrow[r,hook]\arrow[d,"\delta_0"']&\mathcal E\arrow[d,"\delta"]\\
\mathcal E_0\otimes\mathcal A\arrow[r,hook]&\mathcal E\otimes\mathcal A.
\end{tikzcd}
\]
（注意，因 $\mathcal A$ 在 $\mathcal O_S$ 上\Term{flat}{平坦}，\Term{morphism}{态射} $\mathcal E_0\otimes\mathcal A\to\mathcal E\otimes\mathcal A$ 为\Term{injectivity}{单射}。）

则 $\delta_0$ 使 $\mathcal E_0$ 成为 $\mathcal A(G)$-\Term{comodule}{余模}，因而定义了 $G$ 在 $\mathcal E_0$ 上的\Term{action}{作用} $\rho_0$（称为 $\rho$ 在 $\mathcal E_0$ 上的\Term{induced-action}{诱导作用}，并称 $\mathcal E_0$ 在 $\rho$ 下\Term{stable}{稳定}）。
\end{SGAresult}
这立即由定义及 11.6 得到。不过应注意，一般而言，\Term{canonical}{典范}映射 $\mathbf W(\mathcal E_0)\to\mathbf W(\mathcal E)$ 并非\Term{monomorphism}{单态射}。

\textbf{注记 11.7.bis。}\footnote{编者注：补入此注记，它推广了 11.7，并将在 11.10.bis 中使用。}设 $G$ 为\Term{flat}{平坦} $S$-\Term{group-scheme}{群概形}，$\mathcal E$ 为\Term{quasi-coherent}{拟凝聚} $G$-$\mathcal O_S$-模。记 $f$ 为\Term{morphism}{态射} $G\to S$，$\delta$ 为 11.6 中定义的\Term{morphism}{态射} $\mathcal E\to f_*f^*(\mathcal E)$。设 $\mathcal E_0$ 为 $\mathcal E$ 的\Term{quasi-coherent}{拟凝聚} $\mathcal O_S$-\Term{submodule}{子模}。因 $f$ \Term{flat}{平坦}，$f_*f^*(\mathcal E_0)$ 是 $f_*f^*(\mathcal E)$ 的 $\mathcal O_S$-\Term{submodule}{子模}，对 $\phi=f\times f$ 同样如此。因此，若 $\delta$ 在 $\mathcal E_0$ 上的\Term{restriction}{限制} $\delta_0$ 经 $f_*f^*(\mathcal E_0)$ \Term{resolution}{分解}，则它使 $\mathcal E_0$ 成为 $G$-$\mathcal O_S$-模。此时称 $\mathcal E_0$ 为 $\mathcal E$ 的 $G$-\Term{stabilizer}{稳定子}模。

\begin{SGAresult}{定义}{11.8.0}
\footnote{编者注：陈述 11.8、11.9 只涉及\Term{coalgebra}{余代数}上的\Term{comodule}{余模}概念，故引入定义 11.8.0，并据此改述了 11.8、11.9。}设 $S$ 为\Term{scheme}{概形}。一个 $\mathcal O_S$-\Term{coalgebra}{余代数}是带有两个 $\mathcal O_S$-\Term{module-morphism}{模态射} $\Delta:\mathcal C\to\mathcal C\otimes\mathcal C$、$\varepsilon:\mathcal C\to\mathcal O_S$ 的 $\mathcal O_S$-模 $\mathcal C$，满足下列两个公理（参见第 I 讲 4.2）：
\SourcePage{416}

（CO 1）$\Delta$ 满足\Term{coassociative}{余结合律}，即下图交换：
\[
\begin{tikzcd}
&\mathcal C\otimes\mathcal C\arrow[dr,"\operatorname{id}\otimes\Delta"]&\\
\mathcal C\arrow[ur,"\Delta"]\arrow[dr,"\Delta"']&&\mathcal C\otimes\mathcal C\otimes\mathcal C\\
&\mathcal C\otimes\mathcal C\arrow[ur,"\Delta\otimes\operatorname{id}"']&
\end{tikzcd}
\]
（CO 2）$\varepsilon$ 为\Term{counit}{余单位}，即下列两个复合均为\Term{identity-morphism}{恒等态射}：
\[
\begin{tikzcd}
\mathcal C\arrow[r,"\Delta"]&\mathcal C\otimes\mathcal C\arrow[r,"\operatorname{id}\otimes\varepsilon"]&\mathcal C\otimes\mathcal O_S\arrow[r,"\sim"]&\mathcal C,\\
\mathcal C\arrow[r,"\Delta"]&\mathcal C\otimes\mathcal C\arrow[r,"\varepsilon\otimes\operatorname{id}"]&\mathcal O_S\otimes\mathcal C\arrow[r,"\sim"]&\mathcal C.
\end{tikzcd}
\]
一个（右）$\mathcal C$-\Term{comodule}{余模}是带有 $\mathcal O_S$-\Term{module-morphism}{模态射} $\mu:\mathcal E\to\mathcal E\otimes\mathcal C$ 的 $\mathcal O_S$-模 $\mathcal E$，该\Term{morphism}{态射}满足 11.6 的公理（CM 1）、（CM 2）。

若 $\mathcal C$（相应地，$\mathcal E$）作为 $\mathcal O_S$-模\Term{quasi-coherent}{拟凝聚}，就称其为\Term{quasi-coherent}{拟凝聚}\Term{coalgebra}{余代数}（相应地，\Term{quasi-coherent}{拟凝聚}\Term{comodule}{余模}）。
\end{SGAresult}
设 $A$ 为\Term{commutative-ring}{交换环}，$C$ 为 $A$-\Term{coalgebra}{余代数}，则 $C^\vee=\operatorname{Hom}_A(C,A)$ 是 $A$-\Term{algebra}{代数}。以 $\operatorname{ev}$ 表示自然的求值映射 $C\otimes_A C^\vee\to A$。

\begin{SGAresult}{\Term{lemma}{引理}}{11.8}
设 $C$ 为 $A$-\Term{coalgebra}{余代数}，$V$ 为 $C$-\Term{comodule}{余模}，$M$ 为 $V$ 的 $A$-\Term{submodule}{子模}。假设 $C$ 是\Term{projective-module}{投射} $A$-模。\footnote{编者注：原文假设 $C$ 为\Term{free-module}{自由} $A$-模；J.-P. Serre 指出的向\Term{projective-module}{投射} $A$-模的推广，原列于注记 11.10.1。此处及 11.9 中纳入了这一推广，并相应展开了证明。}设 $c(M)$ 为 $A$-\Term{module-morphism}{模态射}
\[
\begin{tikzcd}[column sep=small]
\theta:M\otimes_A C^\vee\arrow[r,"\mu\otimes\operatorname{id}"]&V\otimes_A C\otimes_A C^\vee\arrow[r,"\operatorname{id}\otimes\operatorname{ev}"]&V
\end{tikzcd}
\]
的像。则 $c(M)$ 是 $V$ 中包含 $M$ 的最小\Term{subcomodule}{子余模}，若 $M$ 为\Term{finite-type}{有限型} $A$-模，$c(M)$ 亦然。称 $c(M)$ 为由 $M$ 生成的\Term{subcomodule}{子余模}。

此外，对每个\Term{ring-morphism}{环态射} $A\to A'$，若 $M'$ 表示 $M\otimes_A A'$ 在 $V'=V\otimes_A A'$ 中的像，则 $c(M')$ 是 $c(M)\otimes_A A'$ 在 $V'$ 中的像，即“$c(M)$ 的形成与\Term{base-change}{基变换}交换”。
\end{SGAresult}
\begin{proof}
首先由（CM 2）有 $M\subset c(M)$；若 $N$ 为 $V$ 中包含 $M$ 的\Term{subcomodule}{子余模}，则 $\mu(M)\subset N\otimes C$，因而 $c(M)\subset N$。

由假设，$C$ 是具有基 $(e_i)_{i\in I}$ 的\Term{free-module}{自由} $A$-模 $L$ 的\Term{direct-summand}{直和因子}。记 $\varphi_i$ 为\Term{linear}{线性}型 $e_i^*$ 在 $C$ 上的\Term{restriction}{限制}，其中 $e_i^*(e_j)=\delta_{ij}$。设 $x\in M$。可写成
\[
\mu(x)=\sum_{i\in J}x_i\otimes e_i,\tag{1}
\]
\SourcePage{417}
其中 $x_i\in V$，$J$ 为 $I$ 的\Term{finite}{有限}\Term{subset}{子集}。于是 $x_i=\theta(x\otimes\varphi_i)$ 属于 $c(Ax)$，故 $c(Ax)=\sum_{i\in J}Ax_i$。因为 $C$ 是 $L$ 的\Term{direct-summand}{直和因子}，写作 $L=C\oplus R$，从而 $V\otimes L=(V\otimes C)\oplus(V\otimes R)$，可得
\[
(c(Ax)\otimes L)\cap(V\otimes C)=c(Ax)\otimes C.
\]
因此，$\mu(x)$ 也可写成
\[
\mu(x)=\sum_{j\in J}x_j\otimes b_j,\qquad b_j\in C.\tag{2}
\]
可写 $\Delta(b_j)=\sum_{i\in I}b_{ij}\otimes e_i$，其中 $b_{ij}\in C$。对（1）应用 $\mu\otimes\operatorname{id}$，对（2）应用 $\operatorname{id}\otimes\Delta$，并使用公理（CM 1），则对每个 $i\in J$ 得
\[
\mu(x_i)=\sum_{j\in J}x_j\otimes b_{ij}\in c(Ax)\otimes C.
\]
这表明 $c(M)$ 是 $V$ 的\Term{subcomodule}{子余模}，因而是 $V$ 中包含 $M$ 的最小\Term{subcomodule}{子余模}。

显然，若 $M$ 是\Term{finite-type}{有限型} $A$-模，$c(M)$ 亦然：若 $M=Ax_1+\cdots+Ax_n$，且 $\mu(x_k)=\sum_i x_{ik}\otimes e_i$，则 $c(M)$ 由 $x_{ik}$ 生成，其中 $k=1,\ldots,n$，$i$ 遍历 $I$ 的某个\Term{finite}{有限}\Term{subset}{子集}。

最后，设 $A\to A'$ 为\Term{ring-morphism}{环态射}，$M'$ 为 $M\otimes A'$ 在 $V'=V\otimes A'$ 中的像。则 $c(M')$（相应地，$c(M)\otimes A'$ 在 $V'$ 中的像）为下图中 $\theta'$（相应地，复合 $\theta'\circ\tau$）的像：
\[
\begin{tikzcd}[column sep=small]
M\otimes A'\otimes C^\vee\arrow[r,"\tau"]&M\otimes\operatorname{Hom}_A(C,A')\arrow[r,"\theta'"]&V'.
\end{tikzcd}
\]
这两个\Term{morphism}{态射}有相同的像。事实上，设 $\psi\in\operatorname{Hom}_A(C,A')$、$x\in M$，置 $\mu(x)=\sum_{i\in J}x_i\otimes e_i$，则
\[
\theta'(x\otimes\psi)=\sum_{i\in J}\psi(e_i)x_i
\]
是 $M\otimes A'\otimes C^\vee$ 的元素 $\sum_{i\in J}x\otimes\psi(e_i)\otimes\varphi_i$ 在 $\theta'\circ\tau$ 下的像。\Term{lemma}{引理}得证。
\end{proof}
此外，还有下述命题。
\begin{SGAresult}{命题}{11.8.bis}
\footnote{编者注：补入了本命题，取自 [Se68]，§1.5，命题 2。}设 $A$ 是\Term{noetherian}{诺特}环，$C$ 为在 $A$ 上\Term{flat}{平坦}的 $A$-\Term{coalgebra}{余代数}，$V$ 为 $C$-\Term{comodule}{余模}，$M$ 为 $V$ 的\Term{finite-type}{有限型} $A$-\Term{submodule}{子模}。则存在 $V$ 的\Term{subcomodule}{子余模} $W$，在 $A$ 上\Term{finite-type}{有限型}且包含 $M$。
\end{SGAresult}
\begin{proof}
因为 $M$ \Term{finite-type}{有限型}，$\Delta_V(M)$ 亦然，故存在 $V$ 的\Term{finite-type}{有限型} $A$-\Term{submodule}{子模} $M'$，使 $\Delta_V(M)\subset M'\otimes_A C$。设 $\pi$ 为\Term{projection}{投影} $V\to V/M'$，$\overline\Delta_V=(\pi\otimes\operatorname{id}_C)\Delta_V$，并置
\[
W=\{x\in V\mid\Delta_V(x)\in M'\otimes_A C\}=\operatorname{Ker}\overline\Delta_V.
\]
它是 $V$ 中包含 $M$ 且含于 $M'$ 的 $A$-\Term{submodule}{子模}（因为 $x=(\operatorname{id}_V\otimes\varepsilon)\Delta_V(x)$），因而在 $A$ 上\Term{finite-type}{有限型}。此外，$(\overline\Delta_V\otimes\operatorname{id}_C)\Delta_V=(\pi\otimes\Delta_C)\Delta_V$
\SourcePage{418}
在 $W$ 上为零，即 $\Delta_V(W)$ 包含在 $\overline\Delta_V\otimes\operatorname{id}_C$ 的核 $K$ 中。但因 $C$ 在 $A$ 上\Term{flat}{平坦}，$K=W\otimes C$，故 $W$ 为 $V$ 的\Term{subcomodule}{子余模}。
\end{proof}

\begin{SGAresult}{\Term{lemma}{引理}}{11.8.1}
\footnote{编者注：补入本\Term{lemma}{引理}；11.9 的证明要用到它。}设 $C$ 为 $A$-\Term{coalgebra}{余代数}，$V$ 为 $C$-\Term{comodule}{余模}，$M$ 为 $V$ 的 $A$-\Term{submodule}{子模}，$f:A\to A'$ 为\Term{faithfully-flat}{忠实平坦}\Term{ring-morphism}{环态射}。假设 $C'=C\otimes A'$ 为\Term{projective-module}{投射} $A'$-模。

（i）则存在 $V$ 中包含 $M$ 的最小\Term{subcomodule}{子余模} $t(M)$，且若 $M$ 为\Term{finite-type}{有限型} $A$-模，$t(M)$ 亦然。此外，“$t(M)$ 的形成与\Term{base-change}{基变换}交换”。

（ii）更准确地说，$C$ 为\Term{projective-module}{投射} $A$-模，且 $t(M)=c(M)$。
\end{SGAresult}
\begin{proof}
（ii）由 [RG71]（参见下述命题 11.8.2），$C$ 为\Term{projective-module}{投射} $A$-模。因此可应用\Term{lemma}{引理} 11.8：$c(M)$ 是 $V$ 中包含 $M$ 的最小\Term{subcomodule}{子余模}，当 $M$ 为\Term{finite-type}{有限型} $A$-模时它也是\Term{finite-type}{有限型}的，且其形成与\Term{base-change}{基变换}交换。

为避免使用后来的结果（[RG71] 晚于 SGA 3），略述（i）的直接证明。因 $A\to A'$ \Term{flat}{平坦}，$M'=M\otimes A'$ 是 $V'=V\otimes A'$ 的 $A'$-\Term{submodule}{子模}；又因 $C'$ 为\Term{projective-module}{投射} $A'$-模，$c(M')$ 是 $V'$ 中包含 $M'$ 的最小\Term{subcomodule}{子余模}。以 $V'\otimes A'$ 及 $A'\otimes V'$ 表示在 $V''=V'\otimes_{A'}(A'\otimes A')$ 上由两次\Term{base-change}{基变换} $A'\rightrightarrows A'\otimes A'$、$a'\mapsto a'\otimes1$ 及 $a'\mapsto1\otimes a'$ 得到的两个 $A'\otimes A'$-\Term{comodule}{余模}结构。$A'$-\Term{comodule}{余模} $V'$ 带有 $A'\otimes A'$-\Term{comodule}{余模}的\Term{isomorphism}{同构}
\[
\varphi:V'\otimes A'\xrightarrow{\sim}A'\otimes V',
\qquad(x\otimes a')\otimes b'\longmapsto b'\otimes(x\otimes a'),
\]
它是一个\Term{descent-datum}{下降资料}，即满足 $\varphi_{31}=\varphi_{32}\circ\varphi_{21}$。

因 $M'=M\otimes A'$，$\varphi$ 将 $M'\otimes A'$ 映到 $A'\otimes M'$，因而将 $c(M'\otimes A')$ 映到 $c(A'\otimes M')$。因 $c(M')$ 的形成与\Term{base-change}{基变换}交换，有 $c(M'\otimes A')=c(M')\otimes A'$、$c(A'\otimes M')=A'\otimes c(M')$。故
\[
\varphi(c(M')\otimes A')=A'\otimes c(M'),
\]
于是 $\varphi$ 赋予 $c(M')$ 一个\Term{descent-datum}{下降资料}。由 $(\mathrm{fpqc})$ \Term{descent}{下降}，存在唯一的 $V$ 的\Term{subcomodule}{子余模} $t(M)$，使 $c(M')=t(M)\otimes A'$；且 $t(M)$ 包含 $M$，因为 $t(M)\otimes A'$ 包含 $M'$。若 $N$ 为 $V$ 中包含 $M$ 的\Term{subcomodule}{子余模}，则 $N$ 包含 $t(M)$，因为 $N\otimes A'$ 包含 $c(M')=t(M)\otimes A'$。故 $t(M)$ 是 $V$ 中包含 $M$ 的最小\Term{subcomodule}{子余模}。

最后，设 $A\to B$ 为\Term{ring-morphism}{环态射}。置 $B'=B\otimes A'$，设 $M_B$ 为 $M\otimes B$ 在 $V_B=V\otimes B$ 中的像，$M'_{B'}$ 为 $M'\otimes_{A'}B'$ 在 $V\otimes B'$ 中的像；则 $M_B\otimes_B B'=M'_{B'}$。一方面，将前面的构造应用于 $C_B$ 及\Term{morphism}{态射} $B\to B'$，得
\[
c(M_B\otimes_B B')=t(M_B)\otimes_B B'=t(M_B)\otimes A'.
\]
另一方面，因为 $c(M')$ 的形成与\Term{base-change}{基变换}交换，$c(M'_{B'})$ 是
\[
c(M')\otimes_{A'}B'=t(M)\otimes B\otimes A'
\]
在 $V'\otimes_{A'}B'=V\otimes B\otimes A'$ 中的像。因此 $t(M_B)$ 是 $t(M)\otimes B$ 在 $V_B$ 中的像。
\end{proof}

\SourcePage{419}
\begin{SGAresult}{命题}{11.8.2（Gruson--Raynaud）}
设 $f:A\to A'$ 为\Term{faithfully-flat}{忠实平坦}\Term{morphism}{态射}，则 $f$“使\Term{projective-object}{投射}性\Term{descent}{下降}”，即若 $M$ 为 $A$-模，而 $M\otimes_A A'$ 为\Term{projective-module}{投射} $A'$-模，则 $M$ 为\Term{projective-module}{投射} $A$-模。
\end{SGAresult}
事实上，由 [RG71]，II 2.5.1，$f$“使 \Term{mittag-leffler}{Mittag-Leffler 条件}\Term{descent}{下降}”，故由该文 II 3.1.3，$f$ 使\Term{projective-object}{投射}性\Term{descent}{下降}。\footnote{编者注：顺便指出，该文比 II 2.5.1 更一般的断言 II 2.5.2 已在 [Gr73] 中修正（这不影响\Term{faithfully-flat}{忠实平坦}\Term{morphism}{态射}的情形）。}

\begin{SGAresult}{命题}{11.9}
\footnote{编者注：如编者注（112）所说，已将陈述改写为关于 $\mathcal O_S$-\Term{coalgebra}{余代数} $\mathcal C$ 的\Term{form}{形式}（而非关于满足推论 11.10 所列假设的 $S$-\Term{group-scheme}{群概形} $G$）。另展开了证明；原文仅说“（……）此命题是\Term{lemma}{引理} 11.8 的推论”。}设 $\mathcal C$ 为 $\mathcal O_S$-\Term{coalgebra}{余代数}，$\mathcal E$ 为 $\mathcal C$-\Term{comodule}{余模}，$\mathcal F$ 为 $\mathcal E$ 的 $\mathcal O_S$-\Term{submodule}{子模}，且三者均\Term{quasi-coherent}{拟凝聚}。假设给定 $S$ 的\Term{affine}{仿射}开\Term{covering}{覆盖} $U_\alpha=\operatorname{Spec}A_\alpha$，并对每个 $\alpha$ 给定\Term{faithfully-flat}{忠实平坦}\Term{ring-morphism}{环态射} $A_\alpha\to A'_\alpha$，使 $\Gamma(U_\alpha,\mathcal C)\otimes_{A_\alpha}A'_\alpha$ 是\Term{projective-module}{投射} $A'_\alpha$-模。\footnote{例如，当 $\mathcal C=\mathcal A(G)$，且 $G$ 为\Term{reductive-group}{约化} $S$-\Term{group-scheme}{群概形}时即如此；这将在第 XXII 讲 5.7.8 中证明。}

则存在 $\mathcal E$ 中包含 $\mathcal F$ 的最小\Term{quasi-coherent}{拟凝聚}\Term{subcomodule}{子余模} $t(\mathcal F)$，且若 $\mathcal F$ 为\Term{finite-type}{有限型} $\mathcal O_S$-模，$t(\mathcal F)$ 亦然。此外，对每个\Term{base-change}{基变换} $S'\to S$，若以 $\mathcal F'$ 表示 $\mathcal F\otimes\mathcal O_{S'}$ 在 $\mathcal E'=\mathcal E\otimes\mathcal O_{S'}$ 中的像，则 $t(\mathcal F')$ 为 $t(\mathcal F)\otimes\mathcal O_{S'}$ 在 $\mathcal E'$ 中的像，即“$t(\mathcal F)$ 的形成与\Term{base-change}{基变换}交换”。
\end{SGAresult}
\begin{proof}
对每个 $\alpha$，与 $A_\alpha$-模 $T_\alpha=t(\Gamma(U_\alpha,\mathcal F))$ 相伴的 $\mathcal O_{U_\alpha}$-模 $\mathcal T_\alpha$，由 11.8.1（i），是 $\mathcal E|_{U_\alpha}$ 中包含 $\mathcal F|_{U_\alpha}$ 的最小\Term{quasi-coherent}{拟凝聚}\Term{subcomodule}{子余模}；若 $\mathcal F|_{U_\alpha}$ 为\Term{finite-type}{有限型} $\mathcal O_{U_\alpha}$-模，$\mathcal T_\alpha$ 亦然。

对每个 $\alpha,\beta$，置 $U_{\alpha\beta}=U_\alpha\cap U_\beta$。因 $t(M)$ 的构造与\Term{base-change}{基变换}交换，对每个 $\alpha,\beta$ 有 $\mathcal O_{U_{\alpha\beta}}$-模的\Term{canonical-isomorphism}{典范同构}
\[
\varphi_{\alpha\beta}:T_\beta\otimes_{A_\beta}\mathcal O_{U_{\alpha\beta}}
\xrightarrow{\sim}T_\alpha\otimes_{A_\alpha}\mathcal O_{U_{\alpha\beta}},
\]
它们满足\Term{cocycle-condition}{余循环条件} $\varphi'_{\alpha\gamma}=\varphi'_{\alpha\beta}\circ\varphi'_{\beta\gamma}$，其中 $\varphi'_{\alpha\gamma}$（等等）表示 $\varphi_{\alpha\gamma}$（等等）在 $U_\alpha\cap U_\beta\cap U_\gamma$ 上的\Term{restriction}{限制}。

因此，$\mathcal T_\alpha$ 可\Term{gluing}{粘合}成 $\mathcal E$ 中包含 $\mathcal F$ 的\Term{quasi-coherent}{拟凝聚}\Term{subcomodule}{子余模} $t(\mathcal F)$。留给读者验证，$t(\mathcal F)$ 为 $\mathcal E$ 中包含 $\mathcal F$ 的最小\Term{quasi-coherent}{拟凝聚}\Term{subcomodule}{子余模}，且其形成与\Term{base-change}{基变换}交换。
\end{proof}

\begin{SGAresult}{定义}{11.9.1}
\footnote{编者注：补入本定义，取自 [RG71] 第 82 页下部。因此命题 11.9 的假设就是\Term{coalgebra}{余代数} $\mathcal C$ 为\Term{locally-projective}{局部投射} $\mathcal O_S$-模；推论 11.10 中采用了这一术语。}设 $S$ 为\Term{scheme}{概形}，$\mathcal P$ 为\Term{quasi-coherent}{拟凝聚} $\mathcal O_S$-模。下列条件\Term{equivalence}{等价}：

（i）对 $S$ 的每个\Term{affine}{仿射}\Term{open-subset}{开集} $U$，$\Gamma(U,\mathcal P)$ 为\Term{projective-module}{投射} $\mathcal O_S(U)$-模。
\SourcePage{420}

（ii）存在 $S$ 的\Term{affine}{仿射}开\Term{covering}{覆盖} $(U_\alpha)$，使每个 $\Gamma(U_\alpha,\mathcal P)$ 为\Term{projective-module}{投射} $\mathcal O_S(U_\alpha)$-模。

（iii）存在 $S$ 的\Term{affine}{仿射}开\Term{covering}{覆盖} $(U_\alpha)$，及\Term{faithfully-flat}{忠实平坦}\Term{ring-morphism}{环态射} $A_\alpha=\mathcal O_S(U_\alpha)\to A'_\alpha$，使对每个 $\alpha$，$\Gamma(U_\alpha,\mathcal P)\otimes_{A_\alpha}A'_\alpha$ 为\Term{projective-module}{投射} $A'_\alpha$-模。
\end{SGAresult}
事实上，显然（i）$\Rightarrow$（ii）$\Rightarrow$（iii）。反之，若（iii）成立，11.8.2 蕴含每个 $\Gamma(U_\alpha,\mathcal P)$ 为\Term{projective-module}{投射} $\mathcal O_S(U_\alpha)$-模，故得（ii）。最后，设（ii）成立，且 $V=\operatorname{Spec}A$ 为任意\Term{affine}{仿射}\Term{open-subset}{开集}；可用\Term{finite}{有限}多个\Term{affine}{仿射}\Term{open-subset}{开集} $V_1,\ldots,V_n$ \Term{covering}{覆盖}它，其中每个 $V_i=\operatorname{Spec}A_i$ 至少包含于一个 $V\cap U_\alpha$，从而 $\Gamma(V_i,\mathcal P)$ 为\Term{projective-module}{投射} $A_i$-模。置 $A'=A_1\times\cdots\times A_n$，则 $A\to A'$ \Term{faithfully-flat}{忠实平坦}，且 $\Gamma(V,\mathcal P)\otimes_A A'$ 为\Term{projective-module}{投射} $A'$-模。因此由 11.8.2，$\Gamma(V,\mathcal P)$ 为\Term{projective-module}{投射} $A$-模。

若这些\Term{equivalence}{等价}条件成立，就称 $\mathcal P$ 为\Term{locally-projective}{局部投射} $\mathcal O_S$-模。

\begin{SGAresult}{推论}{11.10}
设 $S$ 为\Term{quasi-compact}{拟紧}且\Term{quasi-separated}{拟分离}的\Term{scheme}{概形}，$G$ 为 $S$-\Term{group-scheme}{群概形}，$\rho$ 为 $G$ 在\Term{quasi-coherent}{拟凝聚} $\mathcal O_S$-模 $\mathcal E$ 上的\Term{linear-action}{线性作用}。假设：

（i）$G$ 满足下列条件之一：

 a）$G$ 在 $S$ 上\Term{affine}{仿射}且\Term{flat}{平坦}；

 b）$G$ 在 $S$ 上\Term{flat}{平坦}、\Term{quasi-compact}{拟紧}、\Term{quasi-separated}{拟分离}，且 $\mathcal E$ \Term{flat}{平坦}；

（ii）$\mathcal A(G)$ 是\Term{locally-projective}{局部投射} $\mathcal O_S$-模。

则 $\mathcal E$ 是其在 $G$ 下\Term{stable}{稳定}的\Term{quasi-coherent}{拟凝聚}\Term{finite-type}{有限型} $\mathcal O_S$-\Term{submodule}{子模}的递增\Term{filtered-family}{滤过族}的\Term{inductive-limit}{归纳极限}。
\end{SGAresult}
由假设（i）及 11.6.1，$\mathcal E$ 带有 $\mathcal A(G)$-\Term{comodule}{余模}结构。另一方面，因 $S$ \Term{quasi-compact}{拟紧}且\Term{quasi-separated}{拟分离}，$\mathcal E$ 是其\Term{quasi-coherent}{拟凝聚}\Term{finite-type}{有限型}\Term{submodule}{子模}的\Term{inductive-limit}{归纳极限}（EGA I，9.4.9 及 EGA IV$_1$，1.7.7）。故将命题 11.9 应用于\Term{coalgebra}{余代数} $\mathcal A(G)$ 即得此推论。

此外，有下列命题。
\begin{SGAresult}{命题}{11.10.bis}
\footnote{编者注：补入本命题，它是 [Th87]，1.4--1.5 的一个特例。该文引用了 Deligne 的一个论证（参见 [Kn71]，III \Term{theorem}{定理} 1.1）；也可注意到它与 11.8.bis 中回顾的 Serre 的论证（[Se68]，命题 2）的相似性。}设 $S$ 为\Term{noetherian}{诺特}\Term{scheme}{概形}，$G$ 为在 $S$ 上\Term{flat}{平坦}、\Term{quasi-compact}{拟紧}、\Term{quasi-separated}{拟分离}的 $S$-\Term{group-scheme}{群概形}，$\mathcal E$ 为\Term{quasi-coherent}{拟凝聚} $G$-$\mathcal O_S$-模，$\mathcal M$ 为 $\mathcal E$ 的\Term{coherent}{凝聚} $\mathcal O_S$-\Term{submodule}{子模}。则 $\mathcal M$ 包含于一个在 $G$ 下\Term{stable}{稳定}的\Term{coherent}{凝聚} $\mathcal O_S$-\Term{submodule}{子模}。
\end{SGAresult}
\begin{proof}
记 $f$ 为\Term{morphism}{态射} $G\to S$，$\tau$ 为\Term{adjunction-morphism}{伴随态射} $\mathcal E\to f_*f^*(\mathcal E)$。由 11.6，$\mathcal E$ 上的 $G$-$\mathcal O_S$-模结构由 $\mathcal O_G$-模 $f^*(\mathcal E)$ 的\Term{automorphism}{自同构} $\theta$ 给出，使\Term{morphism}{态射} $\delta=\theta\circ\tau$ 满足 11.6 的条件（1）、（2）。（[Th87] 中考虑的情形更一般，
\SourcePage{421}
因为该文考虑的是 $G$-\Term{scheme}{概形} $X$ 及 $G$-\Term{equivariant}{等变} $\mathcal O_X$-模 $\mathcal E$，参见第 I 讲第 6 节；这里 $X=S$，带有 $G$ 的\Term{trivial-action}{平凡作用}。）

因 $S$ 是\Term{noetherian}{诺特}的，$\mathcal E$ 是其\Term{coherent}{凝聚}\Term{submodule}{子模} $\mathcal F_\alpha$ 的\Term{filtered}{滤过}\Term{inductive-limit}{归纳极限}（参见 EGA I，9.4.9）。于是 $f^*(\mathcal E)$ 是 $f^*(\mathcal F_\alpha)$ 的\Term{filtered}{滤过}\Term{inductive-limit}{归纳极限}；因 $f$ \Term{flat}{平坦}，后者都是 $f^*(\mathcal E)$ 的\Term{submodule}{子模}。又因 $f$ \Term{quasi-compact}{拟紧}且\Term{quasi-separated}{拟分离}，由 11.0，$\mathcal O_S$-模 $f_*f^*(\mathcal E)$ \Term{quasi-coherent}{拟凝聚}，且是\Term{quasi-coherent}{拟凝聚}\Term{submodule}{子模} $f_*f^*(\mathcal F_\alpha)$ 的\Term{filtered}{滤过}\Term{inductive-limit}{归纳极限}。因此 $\mathcal E$ 是\Term{quasi-coherent}{拟凝聚} $\mathcal O_S$-\Term{submodule}{子模} $\mathcal E_\alpha$ 的\Term{filtered}{滤过}\Term{inductive-limit}{归纳极限}，其中 $\mathcal E_\alpha$ 表示 $f_*f^*(\mathcal F_\alpha)$ 在 $\delta$ 下的\Term{inverse-image}{逆像}，即\Term{composition-morphism}{复合态射}
\[
\overline\delta_\alpha:\mathcal E\xrightarrow{\delta}f_*f^*(\mathcal E)
\longrightarrow f_*f^*(\mathcal E/\mathcal F_\alpha)
\]
的核。因 $\mathcal M$ \Term{coherent}{凝聚}且 $S$ \Term{noetherian}{诺特}，$\mathcal M$ 的任何递增\Term{submodule}{子模}链都\Term{stable}{稳定}，故 $\mathcal M$ 包含于某个 $\mathcal E_\alpha$。下面证明每个 $\mathcal E_\alpha$ 都\Term{coherent}{凝聚}且 $G$-\Term{stable}{稳定}。

设 $u:f^*(\mathcal E)\to\varepsilon_*(\mathcal E)$ 为通过\Term{adjunction}{伴随}对应于 $\mathcal E\to f_*\varepsilon_*(\mathcal E)=\mathcal E$ 的\Term{identity-morphism}{恒等态射}的\Term{morphism}{态射}，则复合
\[
\mathcal E\xrightarrow{\tau}f_*f^*(\mathcal E)
\xrightarrow{f_*(u)}f_*\varepsilon_*(\mathcal E)=\mathcal E
\]
为\Term{identity-morphism}{恒等态射}（参见 11.6（2））。由于有\Term{commutative-diagram}{交换图}
\[
\begin{tikzcd}
\mathcal E_\alpha\arrow[r,"\delta"]\arrow[d,"i_\alpha"']&f_*f^*(\mathcal F_\alpha)\arrow[r,"f_*(u)"]\arrow[d]&\mathcal F_\alpha\arrow[d,"j_\alpha"]\\
\mathcal E\arrow[r,"\tau"']&f_*f^*(\mathcal E)\arrow[r,"f_*(u)"']&\mathcal E,
\end{tikzcd}
\]
其中 $i_\alpha$（相应地，$j_\alpha$）为 $\mathcal E_\alpha$（相应地，$\mathcal F_\alpha$）到 $\mathcal E$ 的包含，故 $i_\alpha$ 经 $j_\alpha$ \Term{resolution}{分解}，即 $\mathcal E_\alpha$ 为 $\mathcal F_\alpha$ 的\Term{submodule}{子模}，因而\Term{coherent}{凝聚}。

最后证明 $\mathcal E_\alpha$ 是 $G$-\Term{stable}{稳定}的（参见 11.7.bis）。用 $h:G\to S$ 表示 $f:G\to S$ 的另一副本，考虑\Term{commutative-diagram}{交换图}
\[
\begin{tikzcd}
G\times_S G\arrow[r,"q"]\arrow[d,"p"']\arrow[dr,"\phi" description]&G\arrow[d,"h"]\\
G\arrow[r,"f"']&S.
\end{tikzcd}
\]
因 $f$ \Term{flat}{平坦}，\Term{exact-sequence}{正合列}
\[
0\longrightarrow\mathcal E_\alpha\longrightarrow\mathcal E
\xrightarrow{\overline\delta_\alpha}h_*h^*(\mathcal E/\mathcal F_\alpha)
\]
给出\Term{exact-sequence}{正合列}
\[
0\longrightarrow f_*f^*(\mathcal E_\alpha)\longrightarrow f_*f^*(\mathcal E)
\xrightarrow{f_*f^*(\overline\delta_\alpha)} f_*f^*h_*h^*(\mathcal E/\mathcal F_\alpha).
\tag{$\dagger$}
\]
\SourcePage{422}
此外，因 $h$ \Term{quasi-compact}{拟紧}且\Term{quasi-separated}{拟分离}，$f$ \Term{flat}{平坦}，\Term{canonical-morphism}{典范态射}
\[
f^*h_*h^*(\mathcal E/\mathcal F_\alpha)
\longrightarrow p_*q^*h^*(\mathcal E/\mathcal F_\alpha)=p_*\phi^*(\mathcal E/\mathcal F_\alpha)
\]
为\Term{isomorphism}{同构}，故（$\dagger$）最右项为 $\phi_*\phi^*(\mathcal E/\mathcal F_\alpha)$。沿用 11.6 的记号，并以 $\pi_\alpha$ 表示\Term{projection}{投影} $\mathcal E\to\mathcal E/\mathcal F_\alpha$，就得到如下\Term{commutative-diagram}{交换图}，其下行\Term{exact}{正合}：
\[
\begin{tikzcd}[column sep=small]
&\mathcal E_\alpha\arrow[r,"\delta"]\arrow[dl,"\delta"']&f_*f^*(\mathcal E)\arrow[d,"f_*(f^*(\pi)\boxtimes\delta_G)"]\\
f_*f^*(\mathcal E_\alpha)\arrow[r]&f_*f^*(\mathcal E)\arrow[r,"f_*f^*(\overline\delta_\alpha)"']&\phi_*\phi^*(\mathcal E/\mathcal F_\alpha).
\end{tikzcd}
\]
因 $f_*(f^*(\pi)\boxtimes\delta_G)\circ\delta$ 在 $\mathcal E_\alpha$ 上为零，可得 $\delta$ 将 $\mathcal E_\alpha$ 映入 $f_*f^*(\mathcal E_\alpha)$，即 $\mathcal E_\alpha$ 为 $G$-\Term{stable}{稳定}。命题得证。
\end{proof}

\textbf{注记 11.10.1。}\footnote{编者注：原注记 11.10.1 的第一部分已并入 11.8、11.9（将“\Term{free-module}{自由}”换为“\Term{projective-module}{投射}”）；下面的反例修正了其第二部分。}在 11.8 中，仅假设 $C$ 是\Term{flat}{平坦} $A$-模并不足够，即使 $A$ 是\Term{principal-ideal-ring}{主理想环}也是如此。事实上，有下列反例，它们由 O. Gabber 及 J.-P. Serre 分别独立地告知我们。

（a）设 $(A,\mathfrak m)$ 为\Term{discrete-valuation-ring}{离散赋值环}，$K$ 为其\Term{fraction-field}{分式域}，$G$ 为 $K$-群 $(\mathbb Z/2\mathbb Z)_K$ 的“\Term{extension-by-zero}{零延拓}”所得的 $A$-\Term{group-scheme}{群概形}。于是\Term{constant-group}{常值群} $(\mathbb Z/2\mathbb Z)_A$，从而其\Term{subgroup}{子群} $G$，通过交换 $v_1,v_2$ \Term{action}{作用}于以 $v_1,v_2$ 为基的\Term{free-module}{自由} $A$-模 $V$。此时 $Av_1$ 不是 $V$ 的 $G$-\Term{submodule}{子模}，但它是 $G$-\Term{submodule}{子模} $Av_1+\mathfrak m^n v_2$（$n\in\mathbb N^*$）的交。因此不存在 $V$ 中包含 $v_1$ 的最小 $G$-\Term{submodule}{子模}。

（b）设 $A$ 为不同于其\Term{fraction-field}{分式域} $K$ 的\Term{integral-domain}{整环}，$G$ 为与下列\Term{Hopf-algebra}{Hopf 代数}对应的\Term{affine}{仿射}\Term{flat}{平坦} $A$-\Term{group-scheme}{群概形}：
\[
\mathcal A(G)=\{P\in K[T]\mid P(0)\in A\},
\]
其中\Term{comultiplication}{余乘法}由 $\Delta(T)=T\otimes1+1\otimes T$ 定义，\Term{counit}{余单位}及\Term{antipode}{对极}分别由 $\varepsilon(T)=0$、$\tau(T)=-T$ 定义。（注意，$G\otimes_A K=\mathbf G_{a,K}$。）

设 $V$ 为\Term{free-module}{自由} $A$-模 $Av_1\oplus Av_2$，$u$ 为由 $u(v_1)=v_2$、$u(v_2)=0$ 定义的 $V$ 的\Term{endomorphism}{自同态}，从而 $u^2=0$。则 $V$ 带有由下式定义的 $\mathcal A(G)$-\Term{comodule}{余模}结构：
\[
\mu(m)=1\otimes m+T\otimes u(m).
\]
$V$ 中包含 $v_1$ 的 $G$-\Term{submodule}{子模}恰好是形如 $Av_1\oplus Iv_2$ 的 $A$-\Term{submodule}{子模}，其中 $I$ 为 $A$ 的非\Term{nilideal}{零理想}；它们的交是 $Av_1$，而 $Av_1$ 不是 $G$-\Term{submodule}{子模}。因此不存在 $V$ 中包含 $v_1$ 的最小 $G$-\Term{submodule}{子模}。（再注意，$C=A\oplus K\cdot T$ 是 $\mathcal A(G)$ 的在 $A$ 上\Term{flat}{平坦}的\Term{subcoalgebra}{子余代数}，且\Term{coaction}{余作用} $\mu:V\to V\otimes\mathcal A(G)$ 经 $V\otimes C$ \Term{resolution}{分解}，所以对这个“非常简单”的\Term{coalgebra}{余代数} $C$ 也得到一个反例。）

最后注意，以上两个例子都是下列构造的特例。设 $A$ 为不同于其\Term{fraction-field}{分式域} $K$ 的\Term{integral-domain}{整环}，$B$ 为在 $A$ 上\Term{free-module}{自由}的 $A$-\Term{Hopf-algebra}{Hopf 代数}。记 $\varepsilon:B\to A$ 为 $B$ 的\Term{augmentation}{增广}，
\SourcePage{423}
$I=\operatorname{Ker}(\varepsilon)$ 为\Term{augmentation-ideal}{增广理想}。因 $B=A\cdot1\oplus I$，容易看出 $B'=\{b\in B\otimes_A K\mid\varepsilon(b)\in A\}$ 是 $B_K$ 的\Term{Hopf-subalgebra}{Hopf 子代数}。若 $V$ 是 $B$-\Term{comodule}{余模}，作为 $A$-模以 $(v_1,\ldots,v_n)$ 为基\Term{free-module}{自由}，且 $\mu_V(v_1)\ne v_1\otimes1$，则 $Av_1$ 不是 $V$ 的\Term{subcomodule}{子余模}，但它是\Term{subcomodule}{子余模} $Av_1+I\cdot V$ 的交，其中 $I$ 遍历 $A$ 的非\Term{nilideal}{零理想}。因此不存在 $V$ 中包含 $v_1$ 的最小\Term{subcomodule}{子余模}。

\begin{SGAresult}{命题}{11.11}
设 $G$ 是域 $k$ 上的\Term{algebraic-group}{代数群}。下列条件\Term{equivalence}{等价}：

（i）$G$ 为\Term{affine}{仿射}的。

（ii）$G$ 为\Term{quasi-affine}{拟仿射}的。

（iii）$G$ \Term{faithful-action}{忠实地作用}于一个\Term{quasi-affine}{拟仿射} $k$-\Term{scheme}{概形} $X$。

（iv）$G$ \Term{linear}{线性}且\Term{faithful-action}{忠实地作用}于一个 $k$-\Term{vector-space}{向量空间}（不必\Term{finite-dimensional}{有限维}）。

（v）$G$ \Term{isomorphism}{同构}于某个群 $\operatorname{GL}(n)_k$ 的\Term{closed-subgroup}{闭子群}。
\end{SGAresult}
\begin{proof}
显然（i）$\Rightarrow$（ii）；而（ii）$\Rightarrow$（iii），因为 $G$ 通过\Term{translation}{平移}\Term{faithful-action}{忠实地作用}于自身。

设 $G$ \Term{faithful-action}{忠实地右作用}于\Term{quasi-affine}{拟仿射} $k$-\Term{scheme}{概形} $X$。\footnote{编者注：以下展开了原文；并选择使 $G$ 在 $X$ 上\Term{right-action}{右作用}，以在 $\mathcal O(X)$ 上得到 $G$ 的\Term{linear}{线性}\Term{left-action}{左作用}。}因 $X$ 为\Term{quasi-affine}{拟仿射}，它是\Term{separated}{分离}且\Term{quasi-compact}{拟紧}的。\footnote{编者注：回顾，称 $k$-\Term{scheme}{概形} $X$ 为\Term{quasi-affine}{拟仿射}，是指它\Term{isomorphism}{同构}于一个\Term{affine}{仿射} $k$-\Term{scheme}{概形}的\Term{quasi-compact}{拟紧}\Term{open-subset}{开集}（EGA II，5.1.1）。}同样，$G$ 是\Term{separated}{分离}的（第 VI$_A$ 讲 0.3），且\Term{quasi-compact}{拟紧}（因其在 $k$ 上\Term{finite-type}{有限型}）。故由 11.1（c），有\Term{canonical-isomorphism}{典范同构}
\[
\begin{aligned}
\mathcal O(X\times G)&=\mathcal O(X)\otimes\mathcal O(G),\\
\mathcal O(X\times G\times G)&=\mathcal O(X)\otimes\mathcal O(G)\otimes\mathcal O(G).
\end{aligned}
\]
由此可得，由\Term{morphism}{态射} $X\times G\to X$ 诱导的\Term{morphism}{态射} $\mu:\mathcal O(X)\to\mathcal O(X)\otimes\mathcal O(G)$ 赋予 $\mathcal O(X)$ 右 $\mathcal O(G)$-\Term{comodule}{余模}结构，即 $G$ \Term{linear}{线性}\Term{left-action}{左作用}于 $k$-\Term{algebra}{代数} $\mathcal O(X)$。因此 $G$ 也\Term{right-action}{右作用}于 $X$ 的\Term{affine-envelope}{仿射包络} $X^{\mathrm{af}}=\operatorname{Spec}\mathcal O(X)$，且\Term{canonical-morphism}{典范态射} $X\to X^{\mathrm{af}}$ 为 $G$-\Term{equivariant}{等变}的。

又因 $X$ 为\Term{quasi-affine}{拟仿射}，$X\to X^{\mathrm{af}}$ 为\Term{open-immersion}{开浸入}（EGA II，5.1.2），故 $G$ 在 $X^{\mathrm{af}}$ 上的\Term{action}{作用}当然也是\Term{faithful-action}{忠实}的。于是 $G$ 在 $k$-\Term{algebra}{代数} $\mathcal O(X)=\mathcal O(X^{\mathrm{af}})$ 上的\Term{linear}{线性}\Term{left-action}{左作用}\Term{faithful-action}{忠实}。这就证明了（iii）$\Rightarrow$（iv）。

现在设 $G$ \Term{faithful-action}{忠实地作用}于 $k$-\Term{vector-space}{向量空间} $V$。由 11.10，$V$ 是在 $G$ 的\Term{action}{作用}下\Term{stable}{稳定}的\Term{finite-dimensional}{有限维}向量子\Term{space}{空间} $V_i$ 的\Term{inductive-limit}{归纳极限}。设 $K_i$ 为 $G$ 在 $V_i$ 上\Term{induced-action}{诱导作用}的核，即\Term{morphism}{态射} $G\to\operatorname{Aut}(V_i)$ 的核，则 $K_i$ 为 $G$ 的闭\Term{subscheme}{子概形}；$G$ 的\Term{action}{作用}\Term{faithful-action}{忠实}，正是说 $K_i$ 的交为 $G$ 的单位\Term{subgroup}{子群}。因 $G$ 是\Term{noetherian}{诺特}的，可知其中某个 $K_i$ 已经只有单位\Term{subgroup}{子群}，故 $G\to\operatorname{Aut}(V_i)$ 是\Term{monomorphism}{单态射}。
\SourcePage{424}
由 1.4.2，它就是\Term{closed-immersion}{闭浸入}，故（iv）$\Rightarrow$（v）。因（v）$\Rightarrow$（i）显然，11.11 得证。
\end{proof}

\textbf{注记 11.11.1。}11.11 可作如下推广。设 $S$ 为\Term{dimension}{维数} $\leq1$ 的\Term{regular}{正则}\Term{locally-noetherian}{局部诺特}\Term{scheme}{概形}，$G$ 为在 $S$ 上\Term{flat}{平坦}、\Term{quasi-compact}{拟紧}、\Term{quasi-separated}{拟分离}的\Term{group-scheme}{群概形}。（此时 $\mathcal A(G)$ 是\Term{torsion-free}{无挠} $\mathcal O_S$-模，因而\Term{flat}{平坦}。）

（a）此时有下列条件的\Term{equivalence}{等价}性：\footnote{编者注：这些条件的\Term{equivalence}{等价}性在补入的第 12 节中证明。}

（i）$G$ 在 $S$ 上\Term{affine}{仿射}。

（ii）$G$ 在 $S$ 上\Term{quasi-affine}{拟仿射}。

（iii）$G$ \Term{faithful-action}{忠实地作用}于一个\Term{quasi-affine}{拟仿射}\Term{flat}{平坦} $S$-\Term{scheme}{概形}。

（iv）$G$ \Term{linear}{线性}且\Term{faithful-action}{忠实地作用}于一个\Term{quasi-coherent}{拟凝聚}\Term{flat}{平坦} $\mathcal O_S$-模。

（b）若 $G$ 还在 $S$ 上\Term{finite-type}{有限型}且 $S$ 是\Term{noetherian}{诺特}的，则这些条件蕴含 $G$ \Term{isomorphism}{同构}于某个 $\operatorname{Aut}(\mathcal E)$ 的\Term{closed-subgroup}{闭子群}，其中 $\mathcal E$ 为\Term{finite-type}{有限型}\Term{locally-free}{局部自由} $\mathcal O_S$-模。\footnote{编者注：这一结论及若干推广将在补入的第 13 节中证明。}

\begin{SGAresult}{\Term{lemma}{引理}}{11.12}
设 $k$ 为域，$G$ 为\Term{affine}{仿射} $k$-\Term{group-scheme}{群概形}。置 $A=\mathcal A(G)$。给定 $x\in A$，存在 $A$ 的\Term{finite-type}{有限型} $k$-\Term{subalgebra}{子代数} $B$，使 $x\in B$，$\Delta(B)\subset B\otimes_k B$，且 $u(B)\subset B$，其中 $u$ 表示与 $G$ 的\Term{inversion}{取逆}\Term{morphism}{态射}对应的 $A$ 的\Term{involution}{对合}。\footnote{编者注：一方面，第 13 节将此推广到 $G$ 为\Term{dimension}{维数} $\leq2$ 的\Term{regular}{正则}底上的\Term{affine}{仿射}\Term{flat}{平坦}群的情形。另一方面，以下对原文作了展开和修正。}
\end{SGAresult}
\begin{proof}
可设 $x\ne0$，于是 $\Delta(x)\ne0$，因为 $(\varepsilon\otimes\operatorname{id})\Delta(x)=x$，其中 $\varepsilon$ 表示 $A$ 的\Term{augmentation}{增广}（\Term{counit}{余单位}）。写成
\[
\Delta(x)=\sum_{j=1}^n e_j\otimes a_j,\qquad n\text{ 最小},\tag{1}
\]
此时 $e_j$（相应地，$a_j$）\Term{linearly-independent}{线性无关}。将 $(e_1,\ldots,e_n)$ 补成 $A$ 的基 $(e_i)_{i\in I}$，对 $j\in J=\{1,\ldots,n\}$ 置
\[
\Delta(e_j)=\sum_{i\in I}e_i\otimes b_{ij}.
\]
对（1）应用 $\Delta\otimes\operatorname{id}$ 及 $\operatorname{id}\otimes\Delta$，由 11.8.0 的公理（CO 1）（另见第 I 讲 4.2 的（HA 1））得等式
\[
\sum_{j\in J}e_j\otimes\Delta(a_j)
=\sum_{\ell\in J}\Delta(e_\ell)\otimes a_\ell
=\sum_{i\in I}e_i\otimes\Bigl(\sum_{\ell\in J}b_{i\ell}\otimes a_\ell\Bigr).
\]
因 $e_i$ \Term{linearly-independent}{线性无关}，可得
\[
\forall j\in J,\qquad\Delta(a_j)=\sum_{\ell\in J}b_{j\ell}\otimes a_\ell.\tag{2}
\]
\SourcePage{425}
令 $B$ 为 $A$ 中由 $b_{ij}$ 及 $u(b_{ij})$（$i,j\in J$）生成的\Term{finite-type}{有限型} $k$-\Term{subalgebra}{子代数}。显然 $u(B)=B$。

对（2）应用 $\Delta\otimes\operatorname{id}$ 及 $\operatorname{id}\otimes\Delta$，再次由（CO 1）得
\[
\sum_{\ell\in J}\Delta(b_{j\ell})\otimes a_\ell
=\sum_{i\in J}b_{ji}\otimes\Delta(a_i)
=\sum_{i,\ell\in J}b_{ji}\otimes b_{i\ell}\otimes a_\ell.
\]
因 $a_\ell$ \Term{linearly-independent}{线性无关}，\Term{pushout}{推出}
\[
\forall j,\ell\in J,\qquad\Delta(b_{j\ell})=\sum_{i\in J}b_{ji}\otimes b_{i\ell}.\tag{3}
\]
因 $\Delta\circ u=(u\times u)\circ v\circ\Delta$，其中 $v(a\otimes b)=b\otimes a$，还有
\[
\forall j,\ell\in J,\qquad\Delta(u(b_{j\ell}))=\sum_{i\in J}u(b_{i\ell})\otimes u(b_{ji}).\tag{4}
\]
因 $\Delta$ 为\Term{algebra}{代数}\Term{homomorphism}{同态}，由（3）、（4）得 $\Delta(B)\subset B\otimes_k B$。最后，11.8.0 的公理（CO 2）（另见第 I 讲 4.2 的（HA 2））表明 $a_j=\sum_{i\in I}\varepsilon(a_j)b_{ij}$，且 $x=\sum_{j\in J}\varepsilon(e_j)a_j$，故 $x\in B$。
\end{proof}

\begin{SGAresult}{命题}{11.13}
设 $k$ 为域，$G$ 为\Term{algebra}{代数}为 $A$ 的\Term{affine}{仿射} $k$-\Term{group-scheme}{群概形}。则 $G$ 是一个\Term{finite-type}{有限型}\Term{affine}{仿射} $k$-\Term{group-scheme}{群概形}的递增\Term{filtered}{滤过}系统的\Term{projective-limit}{投射极限}，其\Term{transition-morphism}{转移态射}均\Term{faithfully-flat}{忠实平坦}。
\end{SGAresult}
\begin{proof}
若 $B,B'$ 为 $A$ 中对 $\Delta,u$ \Term{stable}{稳定}的两个\Term{finite-type}{有限型}\Term{subalgebra}{子代数}，则由 $B,B'$ 生成的\Term{subalgebra}{子代数}亦然。因此由\Term{lemma}{引理} 11.12，$A$ 是一个对 $\Delta,u$ \Term{stable}{稳定}的\Term{finite-type}{有限型}\Term{subalgebra}{子代数}的递增\Term{filtered-family}{滤过族} $(B_i)_{i\in I}$ 的\Term{inductive-limit}{归纳极限}。此时，每个 $B_i$ 连同 $u$ 的\Term{restriction}{限制}和由 $\Delta$ 得到的\Term{morphism}{态射} $B_i\to B_i\otimes_k B_i$，都是\Term{Hopf-algebra}{Hopf 代数}，故由第 I 讲 4.2，它是某个在 $k$ 上\Term{finite-type}{有限型}的\Term{affine}{仿射} $k$-\Term{group-scheme}{群概形} $G_i$ 的\Term{algebra}{代数}。最后，因 $A=\varinjlim B_i$，有 $G=\varprojlim G_i$（参见 EGA IV$_3$，8.2.3）。由下述\Term{lemma}{引理}，\Term{transition-morphism}{转移态射}均\Term{faithfully-flat}{忠实平坦}。
\end{proof}
\begin{SGAresult}{\Term{lemma}{引理}}{11.14}
设 $k$ 为域，$u:G\to H$ 为\Term{finite-type}{有限型}\Term{affine}{仿射} $k$-\Term{group-scheme}{群概形}之间的\Term{morphism}{态射}，$u^\natural:B\to A$\footnote{编者注：把对应于 $u:\operatorname{Spec}(A)\to\operatorname{Spec}(B)$ 的\Term{morphism}{态射} $B\to A$ 记为 $u^\natural$（而非 $u^\circ$）。}为对应的 $k$-\Term{algebra}{代数}\Term{morphism}{态射}。则 $u$ \Term{faithfully-flat}{忠实平坦}的充要条件是 $u^\natural$ \Term{injectivity}{单射}。
\end{SGAresult}
\begin{proof}
必要性显然（参见 EGA $0_{\mathrm I}$，6.6.1）。证明充分性。置 $N=\operatorname{Ker}u$。由第 VI$_A$ 讲 3.3.2、5.4.1，$G/N$ 是\Term{finite-type}{有限型} $k$-\Term{group-scheme}{群概形}，$u$ \Term{resolution}{分解}成
\[
G\xrightarrow{p}G/N\xrightarrow{v}H,
\]
其中 $p$ \Term{faithfully-flat}{忠实平坦}，$v$ 为\Term{closed-immersion}{闭浸入}。因 $H$ 为\Term{affine}{仿射}\Term{scheme}{概形}，$G/N$ 为\Term{affine}{仿射}\Term{scheme}{概形}，且\Term{morphism}{态射} $v^\natural:B\to\mathcal O(G/N)$ \Term{surjectivity}{满射}（参见 EGA I，4.2.3）。由假设 $u^\natural$ \Term{injectivity}{单射}，而 $u^\natural=p^\natural\circ v^\natural$，故 $v^\natural$ 也\Term{injectivity}{单射}，从而是\Term{isomorphism}{同构}，$v$ 亦然。因 $p$ \Term{faithfully-flat}{忠实平坦}，$u$ 亦然。
\end{proof}

\SourcePage{426}
\begin{SGAresult}{定义}{11.15}
\footnote{编者注：补入了 $G$ \Term{quasi-compact}{拟紧}的假设，并展开了条件（i）、（ii）的\Term{equivalence}{等价}性。}设 $k$ 为域，$G$ 为\Term{quasi-compact}{拟紧} $k$-\Term{group-scheme}{群概形}，$V$ 为带有 $G$ 的 $k$-\Term{linear-action}{线性作用}的 $k$-\Term{vector-space}{向量空间}，因而由 11.6.1 带有 $\mathcal A(G)$-\Term{comodule}{余模}结构 $\delta:V\to V\otimes\mathcal A(G)$。设 $v\in V$ 非零。下列条件\Term{equivalence}{等价}：

（i）存在 $\lambda\in\mathcal A(G)$（必然唯一），使 $\delta(v)=v\otimes\lambda$。

（ii）对每个 $k$-\Term{algebra}{代数} $R$ 及每个 $g\in G(R)$，有 $g\cdot v\in Rv$（即存在 $f\in R$，必然唯一，使在 $V\otimes R$ 中有等式 $g\cdot v=v\otimes f$）。

事实上，显然（i）$\Rightarrow$（ii）。反之，若（ii）成立，将其应用于 $R=\mathcal A(G)$ 及 $g=\operatorname{id}_{\mathcal A(G)}$，便得存在唯一的 $\lambda\in\mathcal A(G)$，使 $\delta(v)=v\otimes\lambda$。

若 $v$ 满足这些条件，就称 $v$ 为 $G$ 下的\Term{semi-invariant}{半不变向量}，称 $\lambda$ 为 $v$ 的\Term{weight}{权}；也说“$v$ 为\Term{weight}{权} $\lambda$ 的\Term{semi-invariant}{半不变量}”。
\end{SGAresult}
记 $\Delta$ 为 $\mathcal A(G)$ 的\Term{comultiplication}{余乘法}，则等式
\[
v\otimes\lambda\otimes\lambda=(\delta\otimes\operatorname{id})(\delta(v))
=(\operatorname{id}\otimes\Delta)(\delta(v))=v\otimes\Delta(\lambda)
\]
蕴含 $\Delta(\lambda)=\lambda\otimes\lambda$。因此 $\lambda$ 定义\Term{Hopf-algebra}{Hopf 代数}\Term{morphism}{态射}
\[
\mathcal A(\mathbf G_{m,k})=k[T,T^{-1}]\longrightarrow\mathcal A(G),\qquad T\longmapsto\lambda,
\]
从而定义 $k$-\Term{group-morphism}{群态射} $\lambda:G\to\mathbf G_{m,k}$，即 $\lambda$ 为 $G$ 的\Term{character}{特征标}，称为与\Term{semi-invariant}{半不变向量} $v$ 相伴的\Term{character}{特征标}。

\begin{SGAresult}{\Term{lemma}{引理}}{11.16.0}
\footnote{编者注：补入本\Term{lemma}{引理}，取自 [DG70]，§II.2，3.5。}设 $k$ 为域，$H$ 为\Term{affine}{仿射} $k$-\Term{group-scheme}{群概形}，$V$ 为 $n$ 维 $H$-模，$U$ 为 $V$ 的 $d$ 维向量子\Term{space}{空间}。考虑直线 $D=\bigwedge^d U\subset\bigwedge^d V$。则 $U$ 在 $H$ 下\Term{stable}{稳定}的充要条件是 $D$ 在 $H$ 下\Term{stable}{稳定}。
\end{SGAresult}
\begin{proof}
必要性显然，证明充分性。可设 $d<n$。设 $(e_1,\ldots,e_d)$ 为 $U$ 的一组基，将其补成 $V$ 的基 $(e_1,\ldots,e_n)$。对每个 $k$-\Term{algebra}{代数} $R$，$V_R=V\otimes R$ 是\Term{free-module}{自由} $R$-模，且
\[
U_R=\{v\in V_R\mid v\wedge(e_1\wedge\cdots\wedge e_d)=0\}
\]
（因为对 $i>d$，$e_i\wedge e_1\wedge\cdots\wedge e_d$ 在 $\bigwedge_R^{d+1}V_R$ 中\Term{linearly-independent}{线性无关}）。因 $H(R)$ 在 $\bigwedge_R^\bullet(V_R)$ 上按
\[
h(x_1\wedge\cdots\wedge x_s)=h(x_1)\wedge\cdots\wedge h(x_s)
\]
\Term{action}{作用}，可知若 $H(R)$ \Term{stable}{稳定} $D_R=Re_1\wedge\cdots\wedge e_d$，它也\Term{stable}{稳定} $U_R$。
\end{proof}

\begin{SGAresult}{\Term{theorem}{定理}}{11.16（Chevalley）}
设 $k$ 为域，$G$ 为\Term{affine}{仿射} $k$-\Term{algebraic-group}{代数群}，$H$ 为 $G$ 的闭\Term{subgroup-scheme}{子群概形}。\footnote{编者注：一方面，回顾 $G$ 的每个\Term{subgroup-scheme}{子群概形}都是闭的（1.4.2）。另一方面，将结果表述为通常的\Term{form}{形式}：“$H$ 是 $G$ 的一个\Term{representation}{表示}中一条直线的\Term{stabilizer}{稳定子}”，同时保留了用 $\mathcal A(G)$ 中一列\Term{semi-invariant}{半不变量} $a_1,\ldots,a_n$ 表述的原有\Term{form}{形式}。}则存在\Term{finite-dimensional}{有限维} $G$-模 $V$ 及 $V$ 的一条直线 $D$，使 $H=\operatorname{Norm}_G(D)$，
\SourcePage{427}
即对每个 $k$-\Term{algebra}{代数} $R$，有
\[
H(R)=\{g\in G(R)\mid g(D_R)=D_R\}.
\]
换言之，存在\Term{finite}{有限}多个元素 $a_1,\ldots,a_n\in\mathcal A(G)$，对于 $H$ 的“右”\Term{action}{作用}，它们都是同一\Term{weight}{权} $\lambda$ 的\Term{semi-invariant}{半不变量}（即对每个 $k$-\Term{algebra}{代数} $R$ 及 $g\in G(R)$、$h\in H(R)$，有 $a_i(gh)=\lambda(h)a_i(g)$），使 $H$ 是 $G$ 中使 $a_i$ 都为\Term{semi-invariant}{半不变量}的最大闭\Term{subgroup-scheme}{子群概形}。
\end{SGAresult}
\begin{proof}
记 $\Delta$（相应地，$\varepsilon$）为 $A=\mathcal A(G)$ 的\Term{comultiplication}{余乘法}（相应地，\Term{augmentation}{增广}）。则 $H=\operatorname{Spec}(A/I)$，其中 $I$ 为 $A$ 的某个\Term{ideal}{理想}，包含于 $\operatorname{Ker}\varepsilon$，且 $\Delta(I)\subset I\otimes A+A\otimes I$。设 $B=A/I$，$\pi$ 为\Term{projection}{投影} $A\to B$。考虑由 $(h\varphi)(g)=\varphi(gh)$ 给出的 $H$ 在 $A$ 上的\Term{left-action}{左作用}；相应的 $B$-\Term{comodule}{余模}结构由下式给出：
\[
\overline\Delta:A\xrightarrow{\Delta}A\otimes A\xrightarrow{\operatorname{id}_A\otimes\pi}A\otimes B.
\]
于是 $I$ 是 $A$ 的 $H$-\Term{submodule}{子模}，因为 $\overline\Delta(I)\subset I\otimes B$。

另一方面，$A$ 为\Term{finite-type}{有限型} $k$-\Term{algebra}{代数}，因而\Term{noetherian}{诺特}，故 $I$ 有\Term{finitely-generated}{有限生成}系 $(x_1,\ldots,x_r)$。由 11.8，$x_i$ 包含于某个在 $k$ 上\Term{finite-dimensional}{有限维}的 $G$-\Term{submodule}{子模} $V$。于是 $W=V\cap I$ 为\Term{finite-dimensional}{有限维} $H$-模，记其\Term{dimension}{维数}为 $d$。因 $V$ 包含全部 $x_i$，$W$ 生成\Term{ideal}{理想} $I$。

置 $E=\bigwedge^d V$，设 $(w_1,\ldots,w_d)$ 为 $W$ 的一组基，$\mathcal B=(e_0,\ldots,e_n)$ 为 $E$ 的一组基，其中含有向量 $e_0=w_1\wedge\cdots\wedge w_d$。$G$ 在 $V$ 上的\Term{action}{作用}\Term{canonical}{典范}地决定了 $G$ 在 $E$ 上的\Term{action}{作用}，因而决定了 $A$-\Term{comodule}{余模}结构 $\rho:E\to E\otimes A$。对 $j=0,\ldots,n$，置 $\rho(e_j)=\sum_{i=1}^n e_i\otimes b_{ij}$；在 11.12 的证明中已看到，此时
\[
\Delta(b_{ij})=\sum_{\ell=0}^n b_{i\ell}\otimes b_{\ell j}.\tag{$*$}
\]
置 $a_i=b_{i0}$，即若 $(e_0^*,\ldots,e_n^*)$ 为 $\mathcal B$ 的\Term{dual-basis}{对偶基}，则 $a_i$ 是“\Term{matrix-coefficients}{矩阵系数}”$g\mapsto e_i^*(ge_0)$。另一方面，$H$ 在 $E$ 上的\Term{action}{作用}对应于
\[
\overline\rho:E\xrightarrow{\rho}E\otimes A\xrightarrow{\operatorname{id}_E\otimes\pi}E\otimes B.
\]
因 $W$ 在 $H$ 下\Term{stable}{稳定}，$e_0$ 是 $H$ 下的\Term{semi-invariant}{半不变量}，故 $\overline\rho(e_0)=e_0\otimes\pi(a_0)$，且对 $i=1,\ldots,n$，$a_i=b_{i0}$ 属于 $I$。代入（$*$），得对 $i=0,\ldots,n$，有 $\overline\Delta(a_i)=a_i\otimes\pi(a_0)$，即 $a_i$ 都是 $H$ 下\Term{weight}{权}为 $\pi(a_0)$ 的\Term{semi-invariant}{半不变量}。（此外，将 $E$ 换成由 $e_0$ 生成的 $G$-\Term{submodule}{子模}后，参见 11.8，可设 $a_i$ \Term{linearly-independent}{线性无关}。）

反之，设 $H'=\operatorname{Spec}B'$，其中 $B'=A/I'$，为 $G$ 的闭\Term{subgroup-scheme}{子群概形}，使每个 $a_i$ 都是其下的\Term{semi-invariant}{半不变量}，\Term{weight}{权}为 $\lambda_i\in B'$（特别地，若 $e_0$ 在 $H'$ 下为\Term{weight}{权} $\lambda_0$ 的\Term{invariants}{不变量}，就属于此情形）。证明 $H'=H$。记 $\pi'$
\SourcePage{428}
为\Term{projection}{投影} $A\to B'$；假设蕴含
\[
a_i\otimes\lambda_i=(\operatorname{id}_A\otimes\pi')\Delta(b_{i0})
=\sum_{\ell=0}^n b_{i\ell}\otimes\pi'(a_\ell),
\]
从而 $\lambda_i=\pi'(a_0)$，且对 $\ell=1,\ldots,n$ 有 $a_\ell\in I'$，故 $e_0$ 是 $H'$ 下的\Term{semi-invariant}{半不变量}。由\Term{lemma}{引理} 11.16.0，这蕴含 $W$ 在 $H'$ 下\Term{stable}{稳定}。因\Term{ideal}{理想} $I$ 由 $W$ 生成，它也在 $H'$ 下\Term{stable}{稳定}，故
\[
\Delta(I)\subset I\otimes A+A\otimes I'.
\]
因 $I\subset\operatorname{Ker}\varepsilon$ 且 $(\varepsilon\otimes\operatorname{id}_A)\circ\Delta=\operatorname{id}_A$，可得 $I\subset I'$，从而 $H'\subset H$。
\end{proof}

\begin{SGAresult}{\Term{lemma}{引理}}{11.17.0}
\footnote{编者注：补入本\Term{lemma}{引理}，取自 [DG70]，§III.3，\Term{theorem}{定理} 5.6 的证明（这里滥用记号，以同一个字母 $E$ 表示 $k$-\Term{vector-space}{向量空间}及 $k$-模\Term{scheme}{概形} $\mathbf W(E)=\operatorname{Spec}S(E^*)$）。}设 $k$ 为\Term{algebraically-closed-field}{代数闭域}，$G$ 为\Term{reduced}{约化}的\Term{affine}{仿射} $k$-\Term{algebraic-group}{代数群}，$V$ 为在 $k$ 上\Term{finite-dimensional}{有限维}的 $G$-模，$v\in V$。设 $E$ 为 $V$ 中由向量 $gv$（$g\in G(k)$）生成的向量子\Term{space}{空间}。则 $E$ 是 $V$ 中包含 $v$ 的最小 $G$-\Term{submodule}{子模}，因而\Term{morphism}{态射} $G\times E\to V$、$(g,x)\mapsto gx$ 经 $E$ \Term{resolution}{分解}。
\end{SGAresult}
\begin{proof}
由 11.8，已知存在 $V$ 中包含 $v$ 的最小 $G$-\Term{submodule}{子模} $U$：若 $\mu:V\to V\otimes\mathcal A(G)$ 表示\Term{comodule}{余模}结构，且写成 $\mu(v)=\sum_{i=1}^n v_i\otimes a_i$，其中 $a_i$ \Term{linearly-independent}{线性无关}，则 $U=\operatorname{Vect}(v_1,\ldots,v_n)$。显然 $U$ 包含 $E$，且\Term{morphism}{态射} $G\times U\to V$ 经 $U$ \Term{resolution}{分解}。

反之，$E$ 在\Term{morphism}{态射} $\mu_v:g\mapsto gv$ 下的\Term{inverse-image}{逆像}为 $G$ 中包含\Term{rational-point}{有理点}的闭集；而这些点在 $G$ 中\Term{dense}{稠密}，因为 $G$ 在 $k$ 上\Term{finite-type}{有限型}（参见 EGA IV$_3$，10.4.8）。故 $\mu_v^{-1}(E)=G$，又因 $G$ \Term{reduced}{约化}，$\mu_v$ 经 $E$ \Term{resolution}{分解}，故 $E=U$。
\end{proof}

\begin{SGAresult}{\Term{theorem}{定理}}{11.17（Chevalley）}
设 $k$ 为域，$G$ 为\Term{affine}{仿射} $k$-\Term{group-scheme}{群概形}（不必\Term{finite-type}{有限型}），$N$ 为 $G$ 的闭\Term{subgroup-scheme}{子群概形}，且在 $G$ 中\Term{normal-subgroup}{正规}；则 $(\mathrm{fpqc})$ \Term{quotient-sheaf}{商层} $G/N$ 可由\Term{affine}{仿射} $k$-\Term{group-scheme}{群概形}表示。\footnote{编者注：不使用第 VI$_A$ 讲结果的另一证明，见 [Ta72]，\Term{theorem}{定理} 5.2（另见注记 11.18.5）。}
\end{SGAresult}
\begin{proof}
先设 $G$ 为\Term{finite-type}{有限型}。由第 VI$_A$ 讲 3.2、5.2，$(\mathrm{fpqc})$ \Term{quotient-sheaf}{商层} $G/N$ 可由 $k$-\Term{group-scheme}{群概形} $Q$ 表示，故只须证明 $Q$ 为\Term{affine}{仿射}。证明分几步进行，先设 $k$ 为\Term{algebraically-closed-field}{代数闭域}。\footnote{编者注：以下依据 [DG70]，§III.3，5.6，展开了原文，并修正了错误断言 $(G/N)_{\mathrm{red}}=G_{\mathrm{red}}/N_{\mathrm{red}}$。}

（a）再设 $G$ \Term{reduced}{约化}且\Term{connected}{连通}，$N$ \Term{reduced}{约化}。由 11.16，存在在 $k$ 上\Term{finite-dimensional}{有限维}的 $G$-模 $V$ 及直线 $D=ke_0$，使 $\operatorname{Norm}_G(D)=N$；特别地，$N$ 通过\Term{character}{特征标} $\chi:N\to\mathbf G_{m,k}$ \Term{action}{作用}于 $D$。

固定 $h\in N(k)$。对每个 $g\in G(k)$，有 $hge_0=g(g^{-1}hg)e_0=\chi(g^{-1}hg)ge_0$，故 $\chi(g^{-1}hg)$ 是 $h$ 的一个\Term{characteristic}{特征}值。因此\Term{continuous-map}{连续映射} $\varphi:G(k)\to k$、
\SourcePage{429}
$g\mapsto\chi(g^{-1}hg)$ 只取\Term{finite}{有限}多个值；又因 $G(k)$ \Term{irreducible}{不可约}（它在 $G$ 中\Term{dense}{稠密}），对每个 $g\in G(k)$ 有 $\varphi(g)=\varphi(e)=\chi(h)$，从而
\[
\chi(g^{-1}hg)=\chi(h),\qquad\forall g\in G(k),\ h\in N(k).
\]
设 $E$ 为 $V$ 中由 $ge_0$（$g\in G(k)$）生成的向量子\Term{space}{空间}；由\Term{lemma}{引理} 11.17.0，它是由 $e_0$ 生成的 $V$ 的 $G$-\Term{submodule}{子模}。

由前述结果，两个\Term{morphism}{态射} $N\times E\to E$、$(h,x)\mapsto hx$ 与 $(h,x)\mapsto\chi(h)x$，在\Term{rational-point}{有理点}集 $(N\times E)(k)=N(k)\times E(k)$ 上相同；此\Term{set}{集合}在 $N\times E$ 中\Term{dense}{稠密}。因 $N\times E$ \Term{reduced}{约化}（且 $E$ \Term{separated}{分离}），这两个\Term{morphism}{态射}相等，故 $N$ 以\Term{homothety}{位似}\Term{action}{作用}于 $E$。于是 $N$ 包含在\Term{morphism}{态射} $\rho:G\to\operatorname{GL}(\operatorname{End}_k(E))$ 的核 $K$ 中，其中对每个 $k$-\Term{algebra}{代数} $R$ 及 $g\in G(R)$、$u\in\operatorname{End}_R(E\otimes R)$，置 $\rho(g)(u)=gug^{-1}$。另一方面，若 $g\in K(R)$，则 $g(Re_0)=Re_0$，故 $g\in N(R)$。这说明 $N=K$。由第 VI$_A$ 讲 5.4.1，\Term{morphism}{态射} $G/N\to\operatorname{GL}(\operatorname{End}_k(E))$ 为\Term{closed-immersion}{闭浸入}，故 $G/N$ 为\Term{affine}{仿射}。

（b）现设 $G,N$ \Term{reduced}{约化}，但不要求 $G$ \Term{connected}{连通}。置 $N_0=N\cap G^0$，则由（a），$G^0/N_0$ 为\Term{affine}{仿射}。另一方面，$NG^0$ 是 $G$ 的\Term{normal-subgroup}{正规子群}，而 $G/NG^0$ 作为\Term{finite}{有限}\Term{constant-group}{常值群} $G/G^0$ 的商（参见第 VI$_A$ 讲 5.5.1），同样是\Term{finite}{有限}\Term{constant-group}{常值群}。因此 $G/N$ 是\Term{morphism}{态射} $G/N\to G/NG^0$ 的各\Term{fiber}{纤维}的\Term{direct-sum}{直和}；这些\Term{fiber}{纤维}都\Term{isomorphism}{同构}于 $NG^0/N$，从而由第 VI$_A$ 讲 5.3.3，都\Term{isomorphism}{同构}于 $G^0/N_0$。故 $G/N$ 为\Term{affine}{仿射}。

（c）设 $G$ \Term{reduced}{约化}，$N$ 任意。\Term{morphism}{态射} $G\times N\to N$、$(g,h)\mapsto ghg^{-1}$ 诱\Term{derived-morphism}{导态射} $(G\times N)_{\mathrm{red}}\to N_{\mathrm{red}}$。但因 $G$ \Term{reduced}{约化}且 $k$ 为\Term{algebraically-closed-field}{代数闭域}，有 $(G\times N)_{\mathrm{red}}=G\times N_{\mathrm{red}}$，故 $N_{\mathrm{red}}$ 是 $G$ 的\Term{normal-subgroup}{正规子群}。（注意，$G$ 不\Term{reduced}{约化}时此结论不成立，参见第 VI$_A$ 讲 0.2。）

故由（a），$G'=G/N_{\mathrm{red}}$ 为\Term{affine}{仿射}。另一方面，由第 VI$_A$ 讲 5.6.1，$N'=N/N_{\mathrm{red}}$ 为\Term{finite}{有限} $k$-\Term{group-scheme}{群概形}，故由第 V 讲\Term{theorem}{定理} 4.1，商 $G/N=G'/N'$ 为\Term{affine}{仿射}。

（d）对任意 $G,N$，由 $G\times N\rightrightarrows G$ 经\Term{base-change}{基变换} $G_{\mathrm{red}}\to G$ 得到的\Term{equivalence-relation}{等价关系}为
\[
G_{\mathrm{red}}\times N'\rightrightarrows G_{\mathrm{red}},\qquad
N'=N\cap G_{\mathrm{red}}.
\]
因底\Term{space}{空间}相同（且这个商是商\Term{ringed-space}{赋环空间}），\Term{morphism}{态射} $G_{\mathrm{red}}/N'\to G/N$ 是\Term{homeomorphism}{同胚}。因 $G_{\mathrm{red}}/N'$ \Term{reduced}{约化}（由于 $p:G_{\mathrm{red}}\to G_{\mathrm{red}}/N'$ \Term{faithfully-flat}{忠实平坦}），可得 $(G/N)_{\mathrm{red}}$ 认同于 $G_{\mathrm{red}}/N'$，后者由（c）为\Term{affine}{仿射}。因 $G/N$ 在 $k$ 上\Term{finite-type}{有限型}（参见第 VI$_A$ 讲 3.3.2），由 EGA I，5.1.10 \Term{pushout}{推出} $G/N$ 为\Term{affine}{仿射}。

最后，对任意 $k$，设 $\overline k$ 为 $k$ 的一个\Term{algebraic-closure}{代数闭包}。由 9.2（v），$(G\otimes_k\overline k)/(N\otimes_k\overline k)$ \Term{isomorphism}{同构}于 $(G/N)\otimes_k\overline k$；前者\Term{affine}{仿射}，故后者亦然，从而由 $(\mathrm{fpqc})$ \Term{descent}{下降}（参见 EGA
\SourcePage{430}
IV$_2$，2.7.1），$G/N$ 亦\Term{affine}{仿射}。这就证明了 $G$ 为\Term{finite-type}{有限型}时的 11.17。为推广到一般情形，需要下列\Term{lemma}{引理}。\footnote{编者注：补入此\Term{lemma}{引理}，取自 [DG70]，§III.3，7.1。}
\end{proof}

\begin{SGAresult}{\Term{lemma}{引理}}{11.17.1}
设 $(C_i\to A_i)_{i\in I}$ 为\Term{ring-morphism}{环态射}的\Term{filtered-inductive-system}{滤过归纳系统}，各\Term{morphism}{态射}均\Term{faithfully-flat}{忠实平坦}。则 $A=\varinjlim A_i$ 在 $C=\varinjlim C_i$ 上\Term{faithfully-flat}{忠实平坦}。
\end{SGAresult}
\begin{proof}
由 [BAC]，§I.3，命题 9，\Term{ring-morphism}{环态射} $B\to B'$ \Term{faithfully-flat}{忠实平坦}的充要条件是它\Term{injectivity}{单射}，且 $B'/B$ 为\Term{flat}{平坦} $B$-模。因每个 $C_i\to A_i$ 均\Term{faithfully-flat}{忠实平坦}，有\Term{exact-sequence}{正合列}
\[
0\longrightarrow C_i\longrightarrow A_i\longrightarrow A_i/C_i\longrightarrow0,
\]
且 $A_i/C_i$ 为\Term{flat}{平坦} $C_i$-模，故 $(A_i/C_i)\otimes_{C_i}C$ 为\Term{flat}{平坦} $C$-模。因\Term{inductive-limit}{归纳极限}\Term{exact}{正合}且与\Term{tensor-product}{张量积}交换，得到\Term{exact-sequence}{正合列}
\[
0\longrightarrow C\longrightarrow A\longrightarrow A/C\longrightarrow0
\]
以及\Term{isomorphism}{同构}
\[
\varinjlim\bigl((A_i/C_i)\otimes_{C_i}C\bigr)=(A/C)\otimes_C C=A/C,
\]
由此\Term{pushout}{推出} $A/C$ 为\Term{flat}{平坦} $C$-模。因此 $A$ 在 $C$ 上\Term{faithfully-flat}{忠实平坦}。
\end{proof}
\begin{proof}[\Term{theorem}{定理} 11.17 的证明（续）]
现回到 11.17 在一般情形下的证明，即不假设 $G$ 为\Term{finite-type}{有限型}。置 $\mathcal A(G)=A$，$\mathcal A(N)=B=A/J$。由 11.13，$A$ 是\Term{finite-type}{有限型}\Term{Hopf-subalgebra}{Hopf 子代数}的递增\Term{filtered-family}{滤过族} $(A_i)_{i\in I}$ 的\Term{inductive-limit}{归纳极限}，故 $G$ 是\Term{affine}{仿射} $k$-\Term{algebraic-group}{代数群} $G_i=\operatorname{Spec}(A_i)$ 的\Term{projective-limit}{投射极限}。记 $\Delta$（相应地，$\tau$）为 $A$ 的\Term{comultiplication}{余乘法}（相应地，\Term{antipode}{对极}），并置 $\Delta^2=(\Delta\otimes\operatorname{id}_A)\circ\Delta$。

对每个 $i$，$B_i=A_i/(J\cap A_i)$ 是 $A_i$ 的商\Term{Hopf-algebra}{Hopf 代数}，故 $N_i=\operatorname{Spec}(B_i)$ 是 $G_i$ 的\Term{closed-subgroup}{闭子群}。又因 $N$ 在 $G$ 中\Term{normal-subgroup}{正规}，由 $(g,n)\mapsto gng^{-1}$ 定义的\Term{morphism}{态射} $G\times N\to G$ 经 $N$ \Term{resolution}{分解}；这\Term{equivalence}{等价}于说二元组 $(A,J)$ 满足下述性质：
\[
\bigl(m_{13}\circ(\Delta\otimes\tau)\circ\Delta\bigr)(J)\subset A\otimes_k J,
\]
其中 $m_{13}$ 表示映射 $a_1\otimes a_2\otimes a_3\mapsto a_1a_3\otimes a_2$。由此可得 $(A_i,A_i\cap J)$ 满足相应性质，故 $N_i$ 在 $G_i$ 中\Term{normal-subgroup}{正规}。另一方面，有 $\varinjlim B_i=B$，故 $\varprojlim N_i=N$。

由前面所得，每个 $(\mathrm{fpqc})$ \Term{quotient-sheaf}{商层} $G_i/N_i$ 可由\Term{affine}{仿射} $k$-\Term{group-scheme}{群概形} $Q_i=\operatorname{Spec}(C_i)$ 表示。置 $C=\varinjlim C_i$。于是得到\Term{affine}{仿射} $k$-\Term{group-scheme}{群概形} $Q_i$ 的\Term{filtered}{滤过}\Term{projective-system}{投射系统}，其\Term{projective-limit}{投射极限} $Q=\varprojlim Q_i$ 为 $k$-\Term{group-scheme}{群概形} $\operatorname{Spec}(C)$（参见 EGA IV$_3$，8.2.3）。于是有 $k$-群\Term{exact-sequence}{正合列}
\[
1\longrightarrow N\longrightarrow G\longrightarrow Q.
\]
证明 $Q$ 表示 $G$ 除以 $N$ 的 $(\mathrm{fpqc})$ \Term{quotient-sheaf}{商层}；为此，只须证明\Term{morphism}{态射} $G\to Q$ 相对于 $(\mathrm{fpqc})$ \Term{topology}{拓扑}为\Term{covering}{覆盖}（参见第 IV 讲 3.3.2.1 及 5.1.7.1）。每个\Term{morphism}{态射} $G_i\to Q_i$ 都\Term{faithfully-flat}{忠实平坦}
\SourcePage{431}
（参见 9.2（xi）），即 $A_i$ 在 $C_i$ 上\Term{faithfully-flat}{忠实平坦}；因为 $A=\varinjlim A_i$、$C=\varinjlim C_i$，由\Term{lemma}{引理} 11.17.1 可得 $A$ 在 $C$ 上\Term{faithfully-flat}{忠实平坦}，故 $G\to Q$ 是\Term{faithfully-flat}{忠实平坦}\Term{morphism}{态射}。此\Term{morphism}{态射}为\Term{affine}{仿射}，故\Term{quasi-compact}{拟紧}，因而相对于 $(\mathrm{fpqc})$ \Term{topology}{拓扑}为\Term{covering}{覆盖}。\Term{theorem}{定理} 11.17 的证明完成。
\end{proof}

\textbf{11.18。补充。}\footnote{编者注：补入了本小节。}此外，从 11.17（及其证明）可\Term{pushout}{推出}下列结果，取自 [DG70]，III §3.7。设 $k$ 为域。先给出下述\Term{lemma}{引理}（参见 [An73]，2.3.3.2），后面将用到它（参见 12.10）。

\begin{SGAresult}{\Term{lemma}{引理}}{11.18.1}
设 $u:G\to G'$ 为 $k$-\Term{group-morphism}{群态射}，$N=\operatorname{Ker}(u)$。假设 $G'$ 为\Term{affine}{仿射}，$G$ 为\Term{finite-type}{有限型}。

（i）\Term{morphism}{态射} $\tau:G/N\to G'$ 是\Term{closed-immersion}{闭浸入}。特别地，$G/N$ 为\Term{affine}{仿射}。

（ii）若\Term{morphism}{态射} $u^\natural:\mathcal O(G')\to\mathcal O(G)$ 还\Term{injectivity}{单射}，则 $\tau$ 为\Term{isomorphism}{同构}。（从而 $G'$ 为\Term{finite-type}{有限型}，且 $u$ \Term{faithfully-flat}{忠实平坦}。）
\end{SGAresult}
\begin{proof}
事实上，已知（11.13）$G'$ 是\Term{affine}{仿射} $k$-\Term{algebraic-group}{代数群} $G'_i$ 的某个\Term{filtered}{滤过}系统的\Term{projective-limit}{投射极限}。记 $u_i$ 为\Term{composition-morphism}{复合态射} $G\to G'\to G'_i$，$N_i$ 为其核。则 $N_i$ 构成 $G$ 的\Term{closed-subgroup}{闭子群}的递减\Term{filtered}{滤过}系统，其交为 $N$。因 $G$ 是\Term{noetherian}{诺特}的，存在指标 $i$，使 $N=N_i$。因 $G,G_i$ 均为\Term{finite-type}{有限型}，由第 VI$_A$ 讲 3.2、5.4.1，商 $G/N$ 是\Term{finite-type}{有限型} $k$-\Term{group-scheme}{群概形}，且 $u_i$ 为\Term{faithfully-flat}{忠实平坦}\Term{projection}{投影} $p:G\to G/N$ 与\Term{closed-immersion}{闭浸入} $\tau_i:G/N\hookrightarrow G_i$ 的复合。考虑\Term{commutative-diagram}{交换图}
\[
\begin{tikzcd}
G\arrow[r,"u"]\arrow[d,"p"']&G'\arrow[d,"q_i"]\\
G/N\arrow[r,"\tau_i"']\arrow[ur,"\tau"]&G_i.
\end{tikzcd}
\]
因 $q_i\circ\tau=\tau_i$ 为\Term{closed-immersion}{闭浸入}，且 $q_i$ \Term{separated}{分离}（$G'$ \Term{separated}{分离}，参见第 VI$_A$ 讲 0.3），$\tau$ 为\Term{closed-immersion}{闭浸入}（参见 EGA I，5.4.4）。故 $G/N$ 为\Term{affine}{仿射}，\Term{morphism}{态射} $\tau^\natural:\mathcal O(G')\to\mathcal O(G/N)$ \Term{surjectivity}{满射}，得（i）。

若 $u^\natural:\mathcal O(G')\to\mathcal O(G)$ 还\Term{injectivity}{单射}，则 $\tau^\natural$ 也\Term{injectivity}{单射}，故 $\tau^\natural$ 为\Term{isomorphism}{同构}，因而 $\tau$ 亦然（因 $G',G/N$ 都\Term{affine}{仿射}）。这就证明了（ii）。
\end{proof}

\begin{SGAresult}{\Term{theorem}{定理}}{11.18.2}
设 $G,G'$ 为两个\Term{affine}{仿射} $k$-\Term{group-scheme}{群概形}，其\Term{algebra}{代数}分别为 $A,A'$；设 $u:G\to G'$ 为 $k$-\Term{group-morphism}{群态射}，$N=\operatorname{Ker}(u)$，$\varphi:A'\to A$ 为 $u$ 诱导的\Term{morphism}{态射}。

（i）若 $\varphi$ \Term{injectivity}{单射}，则 $u$ \Term{faithfully-flat}{忠实平坦}，并将 $G'$ 认同于 $G/N$。
\SourcePage{432}

（ii）有 $G/N=\operatorname{Spec}(B)$，其中 $B=\varphi(A')$；故 $u$ 为对应于包含 $B\hookrightarrow A$ 的\Term{faithfully-flat}{忠实平坦}\Term{morphism}{态射} $G\to G/N$，与对应于\Term{surjectivity}{满射} $A'\to B$ 的\Term{closed-immersion}{闭浸入} $G/N\hookrightarrow G'$ 的复合。此外，$N$ 在 $G$ 中由\Term{ideal}{理想} $AB_+$ 定义，其中 $B_+$ 表示 $B$ 的\Term{augmentation-ideal}{增广理想}。

（iii）特别地，若 $u$ 为\Term{monomorphism}{单态射}，则它为\Term{closed-immersion}{闭浸入}。
\end{SGAresult}
\begin{proof}
（i）设 $\varphi$ \Term{injectivity}{单射}，并将 $A'$ 认同于 $A$ 的一个\Term{Hopf-subalgebra}{Hopf 子代数}。由 11.13，$A$ 是\Term{Hopf-subalgebra}{Hopf 子代数} $A_i=\mathcal O(G_i)$ 的\Term{filtered}{滤过}并，其中各 $A_i$ 在 $k$ 上\Term{finite-type}{有限型}；记 $G'_i=\operatorname{Spec}(A'_i)$，其中 $A'_i=A'\cap A_i$，并以 $N_i$ 表示由包含 $A'_i\hookrightarrow A_i$ 诱导的\Term{morphism}{态射} $G_i\to G'_i$ 的核。由上一\Term{lemma}{引理}，有 $G_i/N_i\simeq G'_i$，故对每个 $i$ 得到\Term{exact-sequence}{正合列}
\[
1\longrightarrow N_i\longrightarrow G_i\xrightarrow{p_i}G'_i\longrightarrow1,
\]
其中 $p_i$ \Term{faithfully-flat}{忠实平坦}。因 $G=\varprojlim_i G_i$，$G'=\varprojlim_i G'_i$，得到\Term{exact-sequence}{正合列}
\[
1\longrightarrow N\longrightarrow G\xrightarrow{p}G/N,
\]
其中置 $N=\varprojlim_i N_i$。此外，由\Term{lemma}{引理} 11.17.1，$p$ \Term{faithfully-flat}{忠实平坦}（且\Term{affine}{仿射}），故 $G'$ 表示 $G$ 除以 $N$ 的 $(\mathrm{fpqc})$ \Term{quotient-sheaf}{商层}。这就证明了（i）。

一般情形中，$B=\varphi(A')$ 是 $A$ 的\Term{Hopf-subalgebra}{Hopf 子代数}；记 $H$ 为 $k$-\Term{group-scheme}{群概形} $\operatorname{Spec}(B)$，$N'$ 为由包含 $B\hookrightarrow A$ 诱导的\Term{morphism}{态射} $G\to H$ 的核。由（i），$H$ 认同于 $G/N'$，故 $u$ 是\Term{projection}{投影} $G\to G/N'$ 与由\Term{surjectivity}{满射} $A'\to B$ 诱导的\Term{closed-immersion}{闭浸入} $G/N'\hookrightarrow G'$ 的复合。因此 $N'=N$。再由 9.2（ii），有\Term{cartesian}{笛卡尔}方块
\[
\begin{tikzcd}
N\arrow[r]\arrow[d]&G\arrow[d,"p"]\\
\operatorname{Spec}(k)\arrow[r,"\varepsilon"']&G',
\end{tikzcd}
\]
其中 $\varepsilon$ 是 $G'$ 的\Term{unit-section}{单位截面}，对应于\Term{augmentation}{增广}\Term{morphism}{态射} $B\to k$。故 $N$ 在 $G$ 中由\Term{ideal}{理想} $AB_+$ 定义。这就证明了（ii），（iii）随之得到。
\end{proof}

\textbf{注记 11.18.3。}设 $G$ 为\Term{affine}{仿射} $k$-\Term{group-scheme}{群概形}，$N$ 为\Term{normal-subgroup}{正规} $k$-\Term{subgroup}{子群}。因\Term{morphism}{态射} $p:G\to G/N$ \Term{faithfully-flat}{忠实平坦}且\Term{quasi-compact}{拟紧}，由第 IV 讲 3.3.3.2，$\mathcal O(G/N)$ 是 $\mathcal O(G)$ 中由右 $N$-\Term{invariant}{不变}函数 $\varphi$ 构成的\Term{subalgebra}{子代数}，即满足 $\varphi(gh)=\varphi(g)$ 的函数，其中 $S$ 为任意 $k$-\Term{scheme}{概形}，$g\in G(S)$、$h\in N(S)$。记 $J$ 为 $A=\mathcal O(G)$ 中定义 $N$ 的\Term{ideal}{理想}，则这\Term{equivalence}{等价}于说 $\Delta(\varphi)-\varphi\otimes1\in\mathcal O(G)\otimes J$，其中 $\Delta$ 为 $A$ 的\Term{comultiplication}{余乘法}。

于是上一\Term{theorem}{定理}可用\Term{Hopf-algebra}{Hopf 代数}改述如下。
\begin{SGAresult}{推论}{11.18.4}
设 $k$ 为域，$A$ 为交换 $k$-\Term{Hopf-algebra}{Hopf 代数}，$G=\operatorname{Spec}(A)$。
\SourcePage{433}

（i）若 $B$ 为 $A$ 的\Term{Hopf-subalgebra}{Hopf 子代数}，则 $A$ 在 $B$ 上\Term{faithfully-flat}{忠实平坦}。

（ii）映射 $N\mapsto\mathcal O(G/N)$ 是 $G$ 的\Term{normal-subgroup}{正规子群}集与 $A$ 的\Term{Hopf-subalgebra}{Hopf 子代数}集之间的\Term{bijection}{双射}；其逆映射为 $B\mapsto\operatorname{Spec}(A/AB_+)$。此外，若 $J$ 为 $A$ 中定义 $N$ 的\Term{ideal}{理想}，则
\[
\mathcal O(G/N)=\{x\in A\mid\Delta(x)-x\otimes1\in A\otimes J\}.
\]
\end{SGAresult}

\textbf{注记 11.18.5。}（a）上一\Term{theorem}{定理}的一个推论是，交换\Term{affine}{仿射} $k$-\Term{group-scheme}{群概形}的\Term{category}{范畴}是\Term{abelian-category}{阿贝尔范畴}。关于这一点以及\Term{affine}{仿射} $k$-\Term{group-scheme}{群概形}的其他结果，参见 [DG70]，§III.3，7.4 至 7.8。

（b）最后指出，M. Takeuchi 给出了结果 11.17 至 11.18.4 的另一证明，参见 [Ta72]，§5；他还加强了上面的 11.18.4（i），证明 $A$ 甚至是\Term{projective-module}{投射} $B$-模，参见 [Ta79]，\Term{theorem}{定理} 5（另见 [MW94]，\Term{theorem}{定理} 3.6）。
